% !TeX spellcheck = en_US
\documentclass[12pt,a4paper]{article}
\usepackage{color}
\oddsidemargin=0cm
\evensidemargin=0cm
\textwidth=16.0cm
\topmargin=-1.0cm
\textheight=23.5cm
\usepackage[OT4]{fontenc}
\usepackage[utf8]{inputenc}
\usepackage[polish]{babel}
\usepackage{amsfonts}
\usepackage{amsmath}
\usepackage{amssymb}
\usepackage{fancybox}
\usepackage{graphicx}
\usepackage{xcolor}
\usepackage{colortbl}
%opening
\title{Summary of Professional Accomplishments}
\author{Remigiusz Durka}

\newcommand{\mm}[1]{\textcolor{blue}{#1}}

\setcounter{tocdepth}{4} % to potrzeba, żeby były 
\setcounter{secnumdepth}{4} % paragrafy
\newcommand{\myparagraph}[1]{\paragraph{#1}\mbox{}\\}

\begin{document}
	
%	\maketitle

%%%%%%%%%%%%%%%%%%%%%%%%%%%%%%%%%%%%%%%%%%%%%%%%%%%%%%%%%%%%%%%%%%%%%%%%%%%%%%%%%
\begin{center}
		~~\\[4cm]
		\LARGE
		{\Huge \textbf{Autoreferat}\par}
		%    {\Huge Autoreferat \par}
		\vspace{4cm}
		%   \includegraphics[scale=0.3]{example-image-a}\\
		{\huge Remigiusz Durka}\par
		\vspace{3cm}
		\par
		\begin{minipage}{.5\linewidth}
	\normalsize\centering
	{\large Instytut Fizyki Teoretycznej\\[5pt]
		\large Wydział Fizyki i Astronomii\\[5pt]
		\large Uniwersytet Wrocławski\\[6pt]
		\large Polska}
\end{minipage}
\vfill
\footnotesize{{\selectlanguage{polish}\today} }
%\thispagestyle{empty}
\end{center}
\newpage
\normalsize	
%%%%%%%%%%%%%%%%%%%%%%%%%%%%%%%%%%%%%%%%%%%%%%%%%%%%%%%%%%%%%%%%%%%%%%%%%%%%%%%%
\newpage
\small
\tableofcontents
\normalsize
\newpage


%%%%%%%%%%%%%%%%%%%%%%%%%%%%%%%%%%%%%%%%%%%%%%%%%%%%%%%%%%%%%%%%%%%%%%%%%%%%%%%
\section{Imię i nazwisko}
%\indent
Remigiusz Durka
%%%%%%%%%%%%%%%%%%%%%%%%%%%%%%%%%%%%%%%%%%%%%%%%%%%%%%%%%%%%%%%%%%%%%%%%%%%%%%%%
\section{Posiadane dyplomy i stopnie naukowe}
\begin{itemize}
	\item 2012: Doktorat z fizyki teoretycznej (październik 2007 - czerwiec 2012)\\
	Instytut Fizyki Teoretycznej, Uniwersytet Wrocławski, Polska\\
	Tytuł: „Deformed BF theory as a theory of gravity and supergravity”\\
	Promotor: prof. Jerzy Kowalski-Glikman
	\item 2007: Tytuł magistra z fizyki teoretycznej (październik 2002 - czerwiec 2007)\\
	Wydział Fizyki i Astronomii, Uniwersytet Wrocławski, Polska\\
	Tytuł: „Niezależna od tła grawitacja perturbacyjna”\\
	Promotor: prof. Jerzy Kowalski-Glikman
\end{itemize}

%%%%%%%%%%%%%%%%%%%%%%%%%%%%%%%%%%%%%%%%%%%%%%%%%%%%%%%%%%%%%%%%%%%%%%%%%%%%

\section{3.	Informacja o dotychczasowym zatrudnieniu w jednostkach naukowych}

\begin{itemize}
	\item Kwiecień 2020 - obecnie: Adiunkt\\
	Instytut Fizyki Teoretycznej, Uniwersytet Wrocławski, Polska
	
	\item Październik 2017 - Marzec 2020: Adiunkt\\
	Instytut Fizyki Teoretycznej, Uniwersytet Wrocławski, Polska
	
	\item Listopad 2013 - Listopad 2016: Postdoc\\
	Instituto de F\'{i}sica, Pontificia Universidad Cat\'{o}lica de Valpara\'{i}so (PUCV), Chile\\
	Grant podoktorski FONDECYT\\ 
	Chilijska Comisi\'{o}n Nacional de Investigaci\'{o}n Cient\'{i}fica y Technol\'{o}gica
	
	\item Październik 2012 - Październik 2013: Adiunkt\\
	Instytut Fizyki Teoretycznej, Uniwersytet Wrocławski, Polska
\end{itemize}

%%%%%%%%%%%%%%%%%%%%%%%%%%%%%%%%%%%%%%%%%%%%%%%%%%%%%%%%%%%%%%%%%%%%%%%%%%%

\section{Omówienie osiągnięć}
%Description of the achievements, set out in art. 219 para 1 point 2 of the Act, set out in art. 219 para 1 point 2 of the Act(series of scientific articles related thematically, pursuant to art. 219 para 1. point 2b of the Act.)

    %%%%%%%%%%%%%%%%%%%%%%%%%%%%%%%%%%%%%%%%%%%%%%%%%%%%%%%%%%%%%%%%%%%%%%%%%%%
\subsection{Tytuł osiągnięcia} 

Modele grawitacji i supergrawitacji w kontekście rozszerzonych symetrii.

\newpage
\subsection{Wykaz artykułów związanych z osiągnięciem}
%%%%%%%%%%%%%%%%%%%%%%%%%%%%%%%%%%%%%%%%%%%%%%%%%%%%%%%%%%%%%%%%%%%%%%%%%%%
\begin{itemize}
	\item{\textbf{[H1]}} Dadhich, N., \underline{Durka, R.}, Merino, N. and Miskovic, O.\\ Dynamical structure of Pure Lovelock gravity\\ Phys. Rev. D, 93(6), p.064009 (2016)
	
	\item{\textbf{[H2]}} Concha, P.K., \underline{Durka, R.}, Merino, N. and Rodríguez, E.K.\\ New family of Maxwell like algebras\\ Phys. Lett. B, 759, 507-512 (2016)
	
	\item{\textbf{[H3]}} \underline{Durka, R.}\\ Resonant algebras and gravity\\ Journal of Physics A: Mathematical and Theoretical, 50(14), p.145202 (2017)
	
	\item{\textbf{[H4]}} \underline{Durka, R.}, Kowalski-Glikman, J.\\ Resonant algebras in Chern-Simons model of topological insulators\\ Phys.Lett.B 795, 516-520 (2019)
	
	\item{\textbf{[H5]}} \underline{Durka, R.} and Grela, K.\\ On the number of possible resonant algebras\\ Journal of Physics A: Mathematical and Theoretical, 53(35), p.355202 (2020)
	
	\item{\textbf{[H6]}} Concha, P., \underline{Durka, R.} and Rodríguez, E.\\Resonant superalgebras and N=1 supergravity theories in three spacetime dimensions\\ Phys. Lett. B, 808, p.135659 (2020)
	
	\item{\textbf{[H7]}} \underline{Durka, R.} and Kowalski-Glikman, J.\\ Corners of gravity: the case of gravity as a constrained BF theory\\ Journal of High Energy Physics, (7), pp.1-24 (2021)
	
	\item{\textbf{[H8]}} \underline{Durka, R.}\\ The first law of black hole thermodynamics for Taub–NUT spacetime\\ International Journal of Modern Physics D, 31(04), p.2250021 (2022)
	
	\item{\textbf{[H9]}} \underline{Durka, R.} and Graczyk, K.M.\\ Resonant superalgebras for supergravity. \\The European Physical Journal C, 82(3) (2022)
	
	\item{\textbf{[H10]}} \underline{Durka, R.} and Graczyk, K.M.\\ N=2 resonant superalgebra for supergravity\\ Phys. Lett. B, 833, p.137366 (2022)
	
	\item{\textbf{[H11]}} \underline{Durka, R.} and Kowalski-Glikman, J.\\ A tale of two theories of gravity in asymptotically Anti-de Sitter spacetime\\ Journal of High Energy Physics (02), 029 (2023)
	
	\item{\textbf{[H12]}} \underline{Durka, R.} and Graczyk, K.M.\\ Resonant enlargements of the Poincar\'{e}/AdS (super)algebras from pattern-based analysis\\ Acta Phys. Polon. Supp. 16 6, 11(2023) 
\end{itemize}	
\newpage

%%%%%%%%%%%%%%%%%%%%%%%%%%%%%%%%%%%%%%%%%%%%%%%%%%%%%%%%%%%%%%%%%
\subsection{Wstęp}

Moje badania koncentrują się na modelach grawitacyjnych oraz leżących u ich podstaw symetriach. Moja praca dotyczyła przede wszystkim algebraicznych podstaw teorii grawitacji i supergrawitacji, obejmując rozszerzenie klasycznych algebr Poincarégo oraz Anti-de Sittera (AdS) dla generatorów symetrii do bardziej ogólnych (rezonansowych) struktur algebraicznych. Ponadto badałem złożone zależności między symetriami a modelami grawitacyjnymi, przeprowadzając szczegółową analizę ich własności algebraicznych, geometrycznych i topologicznych. Analiza ta obejmowała rolę brzegu, symetrii asymptotycznych, zachowanych ładunków oraz ich implikacje dla rozwiązań i teorii grawitacyjnych. Poprzez badanie struktury leżącej u podstaw sformułowań grawitacyjnych pogłębiłem zrozumienie, w jaki sposób człony topologiczne wpływają na obserwable fizyczne. W szczególności byłem w stanie opisać rozróżnienie między symetriami objętościowymi a brzegowymi oraz ich rolę w określaniu fizycznej i matematycznej spójności teorii grawitacji.

%%%%%%%%%%%%%%%%%%%%%%%%%%%%%%%%%%%%%%%%%%%%%%%%%%%%%%%%%%%%%%%%%%
\subsubsection{Od Ogólnej Teorii Względności do supergrawitacji}

Ogólna Teoria Względności (OTW) była przedmiotem wielu fundamentalnych rozszerzeń i uogólnień, z wyjątkowo interesującym przykładem supergrawitacji. Ta niezwykła realizacja ucieleśnia ideę supersymetrii między bozonami i fermionami. Supersymetria łączy tutaj bozonowe pole grawitonu o spinie-$2$ z fermionowym polem grawitina o spinie-$3/2$ (Rarita-Schwinger), tworząc tak zwaną supergrawitację (SUGRA) \cite{Freedman:1976xh, VanNieuwenhuizen:1981ae, Townsend:1977qa}.

%$$
%\begin{array}{cc|cc}
%	\textbf{Bosons} &  &  &\textbf{Fermions}\\
%	\textit{Integer Spin} &  &  & \textit{Half-Integer Spin}\\
%	\hline
%	\text{Spin-2} &  &  &\text{Spin-3/2} \\
%	\hline
%	\text{Graviton}\; g_{\mu\nu} &  &  & \text{Gravitino}\;\psi
%\end{array}
%$$
W pośrednim kroku, Ogólna Teoria Względności została wyrażona jako teoria pierwszego rzędu \cite{Hassaine:2016amq}, w której fundamentalnym polem nie jest tensor metryczny, ale zestaw jednoform pól cechowania: vierbein $e^a$ i koneksja spinowa $\omega^{ab}$. Pola te są powiązane z generatorami symetrii grupy Poincar\'e, w szczególności z translacjami czasoprzestrzennymi $P_a$ i transformacjami Lorentza $J_{ab}$. Nowe działanie ma postać:
\begin{align}\label{Einstein}
	I^{4D} &=\frac{1}{32\pi G}\int \mathcal{R}^{ab} (\omega)\wedge e^{c}\wedge e^{d}\,\epsilon_{a b c d} \,,
\end{align}
gdzie $\mathcal{R}^{ab}(\omega) = d\omega^{ab} + \omega^a{}_c \wedge \omega^{cb}$ jest 2-formą krzywizny dla koneksji spinowej.

Podejście to reinterpretuje OTW jako teorię cechowania, podkreślając jej głębokie powiązanie z symetriami i pozwalając na umieszczenie jej na podobnej stopie jak inne fundamentalne oddziaływania w fizyce. Tensor metryczny $g_{\mu\nu}$ w OTW nie jest polem pierwotnym, lecz jest konstruowany z vierbeinu \cite{Durka:2012wd} jako $g_{\mu\nu} = e^a_\mu e^b_\nu \eta_{ab}$. Prowadzi to do bardziej geometrycznej interpretacji grawitacji jako teorii cechowania.

Użycie form różniczkowych, takich jak $\omega^{ab} = \omega^{ab}_\mu dx^\mu$, upraszcza formalizm matematyczny, ujmując wiele ukrytej złożoności w zwięzłym języku geometryczny. Pozwala to na prostsze uogólnienia na wyższe wymiary lub modyfikacje teorii.

Sformułowanie pierwszego rzędu w naturalny sposób pozwala na sprzężenie grawitacji z polami spinorowymi jak fermiony. Spinory transformują się zgodnie z lokalną symetrią Lorentza, a ich pochodna kowariantna wymaga koneksji spinowej $\omega^{ab}$.

W supersymetrycznym rozszerzeniu grawitacji, które obejmuje specyficzny fermionowy generator supersymetrii $Q_{\alpha}$ oraz powiązane z nim pole cechowania, pole grawitina Rarity-Schwingera $\psi^{\alpha}$ o spinie-$\frac{3}{2}$ rozszerza dalej Ogólną Teorię Względności. Pole cechowania w tym ujęciu
\begin{align}
	\mathbb{A} = A^I X_I = \frac{1}{2} \omega^{ab} J_{ab} + \frac{1}{\ell} e^a P_a + \bar{\psi}^{\alpha} Q_{\alpha}\,,
\end{align}
zawiera
\begin{itemize}
	\item $\omega^{ab}$ koneksję spinową związaną z transformacjami Lorentza generowanymi przez $J_{ab}$,
	\item $e^a$ vierbein odpowiadający translacjom czasoprzestrzennym generowanym przez $P_a$,
	\item $\psi^{\alpha}$ pole Rarity-Schwingera związane z fermionowym generatorem superładunku $Q_{\alpha}$.
\end{itemize}

Symetria Lorentza stanowi ważny element fizyki relatywistycznej. Opisuje ona niezmienniczość praw fizycznych pod wpływem obrotów i pchnięć (boostów) w czasoprzestrzeni. Jest ona generowana przez generatory Lorentza $J_{ab}$, które spełniają relacje komutacji algebry Lorentza:
\begin{align}
[J_{ab}, J_{cd}] = \eta_{bc} J_{ad} - \eta_{ac} J_{bd} - \eta_{bd} J_{ac} + \eta_{ad} J_{bc},
\end{align}
gdzie $\eta_{ab}$ jest metryką Minkowskiego.

Translacje w czasoprzestrzeni są generowane przez $P_a$ i komutują między sobą:
\begin{align}
	[P_a, P_b] = 0.
\end{align}
W obecności symetrii Lorentza, generatory $P_a$ i $J_{ab}$ wspólnie tworzą algebrę Poincarégo, której dodatkowe komutatory obejmują:
$$
[J_{ab}, P_c] = \eta_{bc} P_a - \eta_{ac} P_b.
$$
Symetria (Anti-)de Sittera rozszerza algebrę Poincarégo poprzez włączenie stałej kosmologicznej $\Lambda=\mp\frac{3}{\ell^2}$. W tym przypadku translacje $P_a$ oraz transformacje Lorentza $J_{ab}$ są częścią zunifikowanej algebry (A)dS z komutatorem:
\begin{align}
	[P_a, P_b] = \pm \frac{1}{\ell^2} J_{ab},
\end{align}
gdzie $\ell$ jest promieniem (A)dS, a znak ($-$) odpowiada czasoprzestrzeni Anty-de Sittera (AdS), zaś ($+$) de-Sittera (dS).

W przypadku supergrawitacji pełna algebra AdS łączy bozonową algebrę AdS z fermionowymi superładunkami, tworząc bogatą strukturę z zarówno bozonowymi, jak i fermionowymi stopniami swobody. Struktura teorii (zawartość pól, formy krzywizny oraz człony w działaniu – a w konsekwencji także równania pola) jest zdefiniowana przez bazowe (super)algebry generatorów. Istnieją dwa ustalone rodzaje domknięcia superalgebry, prowadzące do dwóch odrębnych działań. Odpowiadają one supersymetrycznym rozszerzeniom Anti-de Sittera z ujemną $\Lambda=-\frac{3}{\ell^2}$, albo Poincaré \cite{Hassaine:2016amq}. Relacje komutacyjne dla generatorów $JP+Q$ w tych dwóch scenariuszach są następujące:
\begin{equation}
	\begin{array}{c|c|c}
		~~ & \text{Poincar\'e Algebra} & \text{AdS Algebra} \\
		\hline
		[J_{ab}, J_{cd}]= & \eta_{b[c} J_{d]a} - \eta_{a[c} J_{d]b} & \eta_{b[c} J_{d]a} - \eta_{a[c} J_{d]b}\\		
		\hline
		[J_{ab}, P_c]= & \eta_{bc} P_a - \eta_{ac} P_b & \eta_{bc} P_a - \eta_{ac} P_b \\
		\hline
		[P_a, P_b]= & 0 & \frac{1}{\ell^2} J_{ab} \\
		\hline
		[J_{ab}, Q_\alpha]= & -\frac{1}{2} (\gamma_{ab} Q)_\alpha & -\frac{1}{2} (\gamma_{ab} Q)_\alpha \\
		\hline
		[P_a, Q_\alpha]= & 0 & -\frac{1}{\ell} (\gamma_a Q)_\alpha \\
		\hline
		\{Q_\alpha, Q_\beta\}= & (\gamma^a C)_{\alpha\beta} P_a & (\gamma^a C)_{\alpha\beta} P_a + \frac{1}{\ell} (\gamma^{ab} C)_{\alpha\beta} J_{ab} \\
	\end{array}
\end{equation}
Poniżej przedstawiono schematycznie \textit{tabele algebr} super-Poincar\'e i super-AdS używane przeze mnie w moich artykułach aby podkreślić ich strukturę, pamiętając jednocześnie o wszystkich współczynnikach, gammach itp.
\begin{equation}\label{table_Poincare}
	Poincare\to
	\begin{array}{c|cc}
		[\,\cdot\,,\,\cdot\,] & J & P \\
		\hline
		J & J & P \\
		P & P & 0 \\
	\end{array}
	\qquad
	\begin{array}{c|cc}
		[\,\cdot\,,\,\cdot\,] & Q & \\
		\hline
		J & Q & \\
		P & 0 & \\
	\end{array}
	\qquad
	\begin{array}{c|c}
		\{\,\cdot\,,\,\cdot\,\} & Q \\
		\hline
		Q & P \\
	\end{array}\quad
\end{equation}

\begin{equation}\label{table_AdS}
	AdS \to\qquad
	\begin{array}{c|cc}
		[\,\cdot\,,\,\cdot\,] & J & P \\
		\hline
		J & J & P \\
		P & P & J \\
	\end{array}
	\qquad\quad
	\begin{array}{c|cc}
		[\,\cdot\,,\,\cdot\,] & Q & \\
		\hline
		J & Q & \\
		P & Q & \\
	\end{array}
	\quad
	\begin{array}{c|c}
		\{\,\cdot\,,\,\cdot\,\} & Q \\
		\hline
		Q & P + J \\
	\end{array}
\end{equation}
Użyta tutaj algebra Liego opisuje komutatory pomiędzy generatorami symetrii: $ [X_a, X_b] = f_{ab}^c X_c $. Tu $X_a$ są generatorami algebry, $f_{ab}^c$ są stałymi struktury, a $[,]$ oznacza komutator. Superalgebry rozszerzają algebry Liego o generatory fermionowe, z antykomutatorem: $\{Y_\alpha, Y_\beta\} = C_{\alpha\beta}^a X_a$.\\

Dla tak zwanej superalgebry Poincar\'e generowanej przez lokalne symetrie cechowania: translacje czasoprzestrzeni $P_a$, transformacje Lorentza $J_{ab}$ i generator supersymetrii $Q_\alpha$, mamy następujące działanie SUGRA
\small
\begin{align}\label{sugra-action}
	I_{\text {ISO}}^{4D} &=\frac{1}{64 \pi G}\int d^4x\, \left(\mathcal{R}_{\mu \nu}^{a b} (\omega)\,e_{\rho}^{c} e_{\sigma}^{d}\right) \epsilon_{a b c d} \epsilon^{\mu \nu \rho \sigma} +\int d^4x\, \left(\frac{1}{2} \bar{\psi}_{\mu} \gamma_{5} \gamma_{a} e_{\nu}^{a} \mathcal{D}_{\rho}^{\omega} \psi_{\sigma}\right) \epsilon^{\mu \nu \rho \sigma}\,,
\end{align}
\normalsize
podczas gdy superalgebra AdS da nam scenariusz ze stałą kosmologiczną $\Lambda=-\frac{3}{\ell^2}$:
\begin{align}\label{sugra-action-ads}
	I_{\text {AdS}}^{4D} &=\frac{1}{64 \pi G}\int d^4x\, \left(\mathcal{R}_{\mu \nu}^{a b} (\omega)\,e_{\rho}^{c} e_{\sigma}^{d}+\frac{1}{\ell^{2}} e_{\mu}^{a} e_{\nu}^{b} e_{\rho}^{c} e_{\sigma}^{d}\right) \epsilon_{a b c d} \epsilon^{\mu \nu \rho \sigma}\notag\\
	&+\int d^4x\, \left(\frac{1}{2} \bar{\psi}_{\mu} \gamma_{5} \gamma_{a} e_{\nu}^{a} \mathcal{D}_{\rho}^{\omega} \psi_{\sigma}+\frac{1}{4 \ell} \bar{\psi}_{\mu} \gamma_{5} \gamma_{a b} e_{\nu}^{a} e_{\rho}^{b} \psi_{\sigma}\right) \epsilon^{\mu \nu \rho \sigma}\,,
\end{align}
Zasada wariacyjna zastosowana do tych działań daje równania pola dla odpowiadających teorii.

%%%%%%%%%%%%%%%%%%%%%%%%%%%%%%%%%%%%%%%%%%%%%%%%%%%%%%%%%%%%%%%%%%%%%%%%%%%%
\subsubsection{Fermionowo $\mathcal{N}$ rozszerzone teorie}

Naturalnym rozwinięciem teorii grawitacji było powiększenie liczby ładunków fermionowych oraz związanych z nimi grawitin, co doprowadziło do powstania tzw. $\mathcal{N}$ supergrawitacji. Temat ten był intensywnie rozwijany przez lata, ujawniając bogatą strukturę możliwych teorii z rosnącą liczbą supersymetrii. Na przykład supergrawitacja $\mathcal{N}=2$ \cite{Gocanin:2019ota} wprowadza dwa fermionowe superładunki i często uwzględnia dodatkowe pola bozonowe, takie jak multiplety wektorowe lub rozszerzenia centralnych ładunków \cite{Haag:1974qh}. Supergrawitacje wyższych rzędów, takie jak $\mathcal{N}=4$ czy $\mathcal{N}=8$, mają głębokie konsekwencje dla teorii unifikacji, ponieważ wykazują dodatkowe symetrie i powiązania z teorią strun. Te osiągnięcia nie tylko pogłębiły nasze zrozumienie fizyki elementarnej, ale również uwydatniły wyzwania związane z utrzymaniem spójności na fundamentalnym poziomie.

%%%%%%%%%%%%%%%%%%%%%%%%%%%%%%%%%%%%%%%%%%%%%%%%%%%%%%%%%%%%%%%%%%%%%%%%%%%%
\subsubsection{Bozonowe rozszerzenia i algebra Maxwella}

Ostatecznie podjęto wysiłki, aby włączyć nowe generatory po stronie bozonowej. Pierwszym takim przykładem jest tzw. algebra Maxwella, jako wynik prac Bacry'ego i Schradera \cite{Bacry:1970ye, Schrader:1972zd}. Dodatkowy bozonowy generator podobny do Lorentza $Z_{ab}$ historycznie wywodzi się z badań nad deformacją symetrii Poincarégo, która uwzględniała obecność pewnego specyficznego stałego pola elektromagnetycznego. Takie podejście doprowadziło do modyfikacji algebry generatorów, w której komutator dwóch translacji równa się nowemu generatorowi podobnemu do Lorentza $Z_{ab}$. Ten typ struktury algebraicznej nazywany jest algebrą Maxwella i ma postać
\begin{align}\label{Maxwell}
	\left[ J_{ab},J_{cd}\right] &=\eta_{bc}J_{ad}-\eta_{ac}J_{bd}+\eta_{ad}J_{bc}-\eta_{bd}J_{ac}\,,\notag\\
	\left[ J_{ab},P_{b}\right] &=\eta _{bc}P_{a}-\eta _{ac}P_{b}\,,\qquad\left[ P_{a},P_{b}\right] =Z_{ab}\,.
\end{align} 
Ostatnio podobna algebra została wprowadzona przez D.~Soroka i V.~Soroka w \cite{Soroka:2006aj, Soroka:2011tc}, ale z bardziej skomplikowaną strukturą,
\begin{align}\label{Soroka}
	\left[ J_{ab},J_{cd}\right] &=\eta_{bc}J_{ad}-\eta_{ac}J_{bd}+\eta_{ad}J_{bc}-\eta_{bd}J_{ac}\,,\notag\\
	\left[ J_{ab},P_{b}\right] &=\eta _{bc}P_{a}-\eta _{ac}P_{b}\,,\qquad\left[ P_{a},P_{b}\right] =Z_{ab}\,,  \notag\\
	\left[Z_{ab},Z_{cd}\right] &=\eta_{bc}Z_{ad}-\eta_{ac}Z_{bd}+\eta_{ad}Z_{bc}-\eta_{bd}Z_{ac}\,,\notag\\
	\left[ Z_{ab},P_{b}\right] &=\eta _{bc}P_{a}-\eta _{ac}P_{b}\,.
\end{align}
Powyższe algebry były przedmiotem wielu badań. Znalazły one interesującą kontynuację w problemie uogólnionych kontrakcji In\"{o}n\"{u}-Wignera \cite{Inonu:1953sp, Lukierski:2010dy}, prowadząc ostatecznie do ciekawego rozwoju znanego jako \textit{ekspansja półgrupowa} (lub w skrócie S-ekspansja) \cite{Izaurieta:2006zz, Andrianopoli:2013ooa, Concha:2016kdz, deAzcarraga:2012zv}. W ramach innego podejścia podobne wyniki (na przykład próby reinterpretacji efektu stałej kosmologicznej jako wynikającego z generatorów $Z_{ab}$ \cite{deAzcarraga:2010sw}) zostały osiągnięte przez kilka innych grup badawczych. Tematyka ta była również częścią mojej pracy doktorskiej \cite{Durka:2011nf, Durka:2011gm}.

%%%%%%%%%%%%%%%%%%%%%%%%%%%%%%%%%%%%%%%%%%%%%%%%%%%%%%%%%%%%%%%%%%%%%%%%%%%%
\subsubsection{Od uogólnionych kontrakcji do S-ekspansji}

Standardowo łączymy przypadek (A)dS (opisany przez stałą kosmologiczną $\Lambda=\mp 3/\ell^2$) z przypadkiem Poincar\'{e}go ($\Lambda=0$), wprowadzając parametr $\ell$, który przeskalowuje generator translacji:
$P_{a} \to \ell \tilde{P}_{a}\,$. Dla przypadku AdS (niezbędnego do sformułowania SUGRA) z $\tilde{J},\tilde{P}$ zapisujemy jawnie $[\ell \tilde{P}_a, \ell \tilde{P}_b]=\tilde{J}_{ab}$. Gdy to równanie zostanie podzielone przez $\ell^2$, a następnie przyjmiemy granicę $\ell \to \infty$, odzyskujemy algebrę Poincar\'{e}go z $[\tilde{P}_a,\tilde{P}_b]=0$.

Jeśli zamiast skalarnych przeskalowań generatory algebry zostaną pomnożone przez odpowiednie elementy $\lambda_i$ pewnej półgrupy (na przykład reprezentowanej przez pewne macierze), otrzymamy:
\begin{align}
	P_{a} \to \lambda_1 \bigotimes \tilde{P}_{a}\qquad J_{ab} \to \lambda_0 \bigotimes \tilde{J}_{ab} \qquad Z_{ab} \to \lambda_2 \bigotimes \tilde{J}_{ab}\,.
\end{align}
W rezultacie otrzymujemy tzw. S-rozszerzoną algebrę \cite{Izaurieta:2006zz}, \cite{Andrianopoli:2013ooa}, oznaczoną jako $S \bigotimes \mathfrak{g}$, gdzie $S=\{\lambda_0,\lambda_1,\lambda_2,\ldots\}$ oraz $\mathfrak{g}=AdS$. W tym scenariuszu jednak nie możemy jednak przyjąć limitu omówionego powyżej.

Zamiast tego musimy rozważyć wynik mnożenia między elementami półgrupy, który ostatecznie przeskalowuje bazowe generatory AdS $\tilde{J}_{ab}, \tilde{P}_{a}$. Znalezienie nowej algebry generatorów sprowadza się zatem do znalezienia odpowiedniej istniejącej półgrupy. Na przykład, relacje komutacyjne dla algebry Soroki-Soroki \cite{Salgado:2013eut} będą określone przez prawo mnożenia, które ustala maksymalną etykietę elementu $i=2$:
\begin{align}\label{product_law}
	\lambda_\alpha \lambda_\beta =
	\begin{cases} 
		\lambda_{\alpha + \beta} & \text{gdy } \alpha + \beta \leq i, \\
		\lambda_{\alpha + \beta - 2 \left\lfloor \frac{n+1}{2} \right\rfloor} & \text{gdy } \alpha + \beta > i,
	\end{cases}
\end{align}
Nowe generatory są wprowadzone przez półgrupowe „skalowania”:
\begin{align}
	J_{ab,(i)}=\lambda_{2i}\tilde{J}_{ab},\quad\text{and}\quad P_{a,(i)}=\lambda_{2i+1}\tilde{P}_{a},\quad 
\end{align}
dla $i=\{0,1,2,...\}$, gdzie
\begin{align}
	J_{ab}&= \lambda_0\tilde{J}_{ab},\quad  P_{a}= \lambda_1\tilde{P}_a,\notag\\
	Z_{ab}&= \lambda_2 \tilde{J}_{ab},\quad U_{a}= \lambda_3\tilde{P}_a, \quad \dots 
\end{align}
dla oryginalnej algebry $\mathfrak{g}=AdS$ określonej przez $\tilde{J}_{ab},\tilde{P}_{a}$ oraz pewne elementy półgrupy $\lambda_i \in S$, które spełniają tzw. \textit{warunek rezonansowy}. Matematycznie, dla części bozonowej, można to interpretować jako wymóg warunku parzystości: Termin „rezonansowości” pochodzi od zastosowania rezonansowego rozkładu elementów półgrupy
\begin{equation}
	\lambda_{even}\cdot\lambda_{even}=\lambda_{even}\,, \qquad \lambda_{even}\cdot\lambda_{odd}=\lambda_{odd}\,, \qquad \lambda_{odd}\cdot\lambda_{odd}=\lambda_{even}\,, \label{Eq:lambdas_rules}  
\end{equation}
który odzwierciedla fakt dekompozycji oryginalnej algebry AdS
\begin{equation}
	[\tilde{J},\tilde{J}] \sim \tilde{J}\,, \qquad [\tilde{J},\tilde{P}] \sim \tilde{P}\,, \qquad [\tilde{P},\tilde{P}] \sim \tilde{J}\,. 
\end{equation}

W odróżnieniu od standardowej kontrakcji In\"{o}n\"{u}-Wignera, nie mamy limitu opisującego przejście od AdS do Poincar\'{e}, gdy stała kosmologiczna zanika przez $\ell \to \infty$. Zamiast tego występuje jedynie wyznaczenie produktów określających pojawienie się nowych generatorów. Ponieważ półgrupa zmienia jedynie typ generatora, produkt $S \times \mathfrak{g}$ zapewnia dziedziczenie stałych struktury, które są budowane według tego samego schematu co początkowa $\mathfrak{g}=AdS$.

Rozszerzenie półgrupowe zapewnia interesujące podejście do systematycznego konstruowania dalszych algebr bozonowych: wystarczy kontynuować dodawanie nowych generatorów, dołączając kolejne $\lambda_3, \lambda_4, \ldots$.

Jednocześnie zero w komutatorach zostało wprowadzone w procesie tzw. redukcji $0_S$ \cite{Izaurieta:2006zz, Salgado:2013eut} lub, jak później zaproponowałem w (\textbf{H3}), poprzez proste uwzględnienie elementu zerowego jako elementu absorbującego w półgrupie.

Dodam, że dla danej tabeli mnożenia półgrupy można natychmiast odczytać odpowiadającą jej algebrę Liego. Rozszerzone generatory można schematycznie zestawić w podobnej tabeli komutacji, oczywiście z uwzględnieniem szczególnych stałych struktury należących do komutatorów $\left[X_{..}, X_{..}\right]$, $\left[X_{..}, X_{.}\right]$ i $\left[X_{.}, X_{.}\right]$. Biorąc jako przykład oryginalną algebrę Maxwella \eqref{Maxwell}, gdzie $0_S\mathbb{T}_M=0$, widzimy, że:
\begin{align}
	\begin{array}{c|ccc}
		\cdot & \lambda_0 & \lambda_1 & \lambda_2 \\
		\hline
		\lambda_0 & \lambda_0 & \lambda_1 & \lambda_2 \\
		\lambda_1 & \lambda_1 & \lambda_2 & 0_S \\
		\lambda_2 & \lambda_2 & 0_S & 0_S
	\end{array}
	\qquad
	\text{semigroup} \leftrightarrow \text{algebra}
	\qquad
	\begin{array}{c|ccc}
		\left[~,~\right] & J_{..} & P_{.} & Z_{..} \\
		\hline
		J_{..} & J_{..} & P_{.} & Z_{..} \\
		P_{.}  & P_{.} & Z_{..} & 0 \\
		Z_{..} & Z_{..} & 0 & 0
	\end{array}
\end{align}
z generatorami zdefiniowanymi jako
\begin{align}
	J_{ab}=\lambda_0\tilde{J}_{ab}\,,\qquad P_{a}=\lambda_1\tilde{P}_a\,,\qquad Z_{ab}=\lambda_2\tilde{J}_{ab}\,.
\end{align}
Relacje komutacyjne mogą być bezpośrednio odczytane z tabeli po prawej stronie lub wyprowadzone jawnie przy użyciu oryginalnej algebry AdS generatorów z tyldą,
\begin{align}
	\left[  J_{ab,}J_{cd}\right]   &  =\lambda_0\lambda_0(\eta_{bc}\tilde{J}_{ad}-\eta_{ac}\tilde{J}_{bd}+\eta
	_{ad}\tilde{J}_{bc}-\eta_{bd}\tilde{J}_{ac})=\eta_{bc}J_{ad}-\eta_{ac}J_{bd}+\eta
	_{ad}J_{bc}-\eta_{bd}J_{ac}\,,\nonumber\\
	\left[  J_{ab},Z_{cd}\right]   &  =\lambda_0\lambda_2(\eta_{bc}\tilde{J}_{ad}-\eta_{ac}\tilde{J}_{bd}+\eta
	_{ad}\tilde{J}_{bc}-\eta_{bd}\tilde{J}_{ac})=\eta_{bc}Z_{ad}-\eta_{ac}Z_{bd}+\eta
	_{ad}Z_{bc}-\eta_{bd}Z_{ac}\,,\nonumber\\
	\left[  Z_{ab,}Z_{cd}\right]   &  =\lambda_2\lambda_2(\eta_{bc}\tilde{J}_{ad}-\eta_{ac}\tilde{J}_{bd}+\eta
	_{ad}\tilde{J}_{bc}-\eta_{bd}\tilde{J}_{ac})=0\,,\nonumber\\
	\left[  J_{ab}%
	,P_{c}\right] & =\lambda_0\lambda_1(\eta_{bc}\tilde{P}_{a}-\eta_{ac}\tilde{P}_{b})=\eta_{bc}P_{a}-\eta_{ac}P_{b}\,,\nonumber\\
	\left[  Z_{ab}%
	,P_{c}\right] & =\lambda_2\lambda_1(\eta_{bc}\tilde{P}_{a}-\eta_{ac}\tilde{P}_{b})=0\,,\nonumber\\
	\left[  P_{a},P_{b}\right]   & =\lambda_1\lambda_1\tilde{J}_{ab}=\lambda_2\tilde{J}_{ab}=Z_{ab}\,.
\end{align}

%%%%%%%%%%%%%%%%%%%%%%%%%%%%%%%%%%%%%%%%%%%%%%%%%%%%%%%%%%%%%%%%%%%%%%%%%%%%
\subsection{Rezonansowe (super)algebry - mój wkład}

Ekspansja półgrupowa dostarczyła spójnego i systematycznego podejścia do odkrywania nowych struktur algebraicznych \textbf{H2} (zob. także \cite{Izaurieta:2006zz, Andrianopoli:2013ooa}), co ostatecznie doprowadziło do rozwoju \textit{algebr rezonansowych} (\textbf{H3, H4, H5}), które wykraczają poza pojedyncze prawo mnożenia półgrupy, takie jak \eqref{product_law}. Jednym z moich wkładów w tej dziedzinie jest włączenie rozszerzeń supersymetrycznych (\textbf{H6, H9, H10, H12}) do struktury rezonansowej, co poszerzyło jej zakres i użyteczność.

W istocie koncepcja rezonansowych (super)algebr stanowi jedno z moich głównych oryginalnych osiągnięć w tej dziedzinie. W kolejnych sekcjach, opierając się w dużej mierze na (\textbf{H6}) i (\textbf{H9}), przedstawię ogólne ramy i kluczowe cechy tych algebr, w tym ich praktyczne zastosowania oraz istotne postępy obliczeniowe. Te osiągnięcia dostarczają systematycznej metodologii odkrywania nowych algebr i dokładnej klasyfikacji ich właściwości.

%%%%%%%%%%%%%%%%%%%%%%%%%%%%%%%%%%%%%%%%%%%%%%%%%%%%%%%%%%%%%%%%%%%%%%%%%%%%
\subsubsection{Konstrukcja rezonansowej superalgebry}\label{Construction}

S-ekspansja nazywana jest \textit{rezonansową}, gdy dekompozycja półgrupy spełnia tę samą strukturę, co podprzestrzenie oryginalnej (super)algebry. W trakcie moich badań odkryłem, że cała struktura rezonansowa może zostać zdefiniowana w inny sposób. W praktyce, aby skonstruować tzw. rezonansowe (super)algebry, należy rozpocząć od tradycyjnej (super)algebry AdS z generatorami $JP(Q)$. Następnie wprowadzamy generatory typu $J$-podobne, $P$-podobne i $Q$-podobne, zachowując ten sam schemat stałych struktury, co dla $JP+Q$.

Struktura \textit{rezonansowa} może być wtedy rozumiana jako wypełnienie wszystkich możliwości w szablonie poniżej w zależności od (rozszerzonego) zestawu generatorów:

\begin{align}\label{empty}
	\left\{
	\begin{aligned} 
		\left[ \framebox(10,10){~~}_{ab},\ \framebox(10,10){~~}_{cd}\right] & =\eta_{bc} \framebox(10,10){~~}_{ad}-\eta_{ac}\ \framebox(10,10){~~}_{bd}-\eta_{bd} \framebox(10,10){~~}_{ac}+\eta_{ad} \framebox(10,10){~~}_{bc}\, ,\\
		\left[\framebox(10,10){~~}_{ab},\ \dashbox(10,10){~~}_{c}\right] & =\eta_{bc} \dashbox(10,10){~~}_{a}-\eta_{ac} \dashbox(10,10){~~}_{b}\, , \\
		\left[\dashbox(10,10){~~}_{a},\ \dashbox(10,10){~~}_{b}\right] & =\framebox(10,10){~~}_{ab}\, ,\\
		\left[ \framebox(10,10){~~}_{ab},\ovalbox{\phantom{I}}_{\alpha }\right]
		&=\frac{1}{2}\,\left( \Gamma _{ab}\right) _{\alpha }^{\beta }\ovalbox{\phantom{I}}_{\beta}\,,\\ 
		\left[ \dashbox(10,10){~~}_{a},\ \ovalbox{\phantom{I}}_{\alpha }\right]&=\frac{1}{2}\,\left( \Gamma _{a}\right) _{\alpha }^{\beta }\ovalbox{\phantom{I}}_{\beta}\,, \\ 
		\left\{ \ovalbox{\phantom{I}}_{~\alpha },\ \ovalbox{\phantom{I}}_{~\beta }\right\} &=-\left(	\Gamma ^{a} C\right) _{\alpha \beta } \dashbox(10,10){~~}_{~a}+\frac{1}{2}\left(\Gamma ^{ab} C\right) _{\alpha \beta } \framebox(10,10){~~}_{~ab}\,.
	\end{aligned} 
	\right.
\end{align}
Dostępny zestaw generatorów zawiera
\begin{align}
	& \framebox(10,10){~~}_{~ab} \to J_{ab}, Z_{ab}, \ldots \quad\; (\rm{Lorentzowsko-podobne})\notag\\
	& \dashbox(10,10){~~}_{~a}\to P_{a}, U_{a}, \ldots \quad\quad\, (\rm{translacyjnie-podobne})\\
	& \ovalbox{\phantom{I}}_{~\alpha } \to Q_\alpha, Y_\alpha, \dots \quad\quad (\rm{superładunkowo-podobne})\notag
\end{align}
i może być dalej rozszerzany, aby uwzględnić jeszcze bardziej rozbudowany zestaw generatorów. Powyżej wprowadzamy macierze $\Gamma_a$, a także $\Gamma_{ab}=1/2(\Gamma_a \Gamma_b-\Gamma_b \Gamma_a)$ oraz macierz sprzężenia ładunkowego $C$, gdzie $a,b,c,d$ są indeksami Lorentza przyjmującymi wartości od $0$ do $3$, a $\alpha,\beta$ są indeksami spinorowymi.

Oczywiście nie wszystkie możliwości są fizycznie wykonalne; dlatego wprowadzamy pewną liczbę wymagań. Dopuszczamy również, aby zero pojawiało się jako wynik komutatorów i antykomutatorów po prawej stronie równań~\eqref{empty}.\\~~\\
%
Algebraiczne i fizyczne wymagania, które muszą zostać spełnione, to:
\begin{itemize}
	\item zachowanie tych samych stałych struktury co w oryginalnej superalgebrze AdS;
	\item zachowanie przez generator Lorentza, tj. dla wszystkich generatorów $[J, X]\sim X$;
	\item antykomutator $\{Q,Q\}$ musi być różny od zera;
	\item spełnienie uogólnionych tożsamości super-Jacobiego.
\end{itemize}
Dwa pierwsze warunki wykluczają jakiekolwiek naruszenie relacji $[J, P] \sim P$. 
Odrzucenie ograniczenia $[J, P]$ prowadziłoby do problemu z definiowaniem torsji. Powoduje to komplikacje w uzyskaniu odpowiednich transformacji vierbeinu i koneksji spinowej oraz w ustanowieniu równoważności z metryczną teorią grawitacji, co jest zagwarantowane przez zależność $\omega(e)$.

Powyższe ograniczenia zapewniają struktury algebraiczne odpowiednie do konstruowania działań, choć niektóre mogą być uznane za niestandardowe lub egzotyczne. Dla poprawnej konstrukcji supergrawitacji wymagamy dodatkowego warunku:
\begin{equation}
	\label{Eq_QQSugra}
	\{Q,Q\}\sim P+\dots
\end{equation}  

Zauważmy, że taka konfiguracja, choć zawiera dodatkowy ładunek fermionowy, nie obejmuje centralnego ładunku \cite{Gocanin:2019ota, Haag:1974qh}. Takie rozszerzenie pozwala na zakodowanie generatora $U(1)$ $T$ w ramach wyniku antykomutacji $Q$ z $Y$. Rozważania te zostały omówione w (\textbf{H10, H12}).

\subsubsection{Przegląd rezonansowych superalgebr}\label{Overview}

Podążając za programem poszukiwania wszystkich możliwości (super)algebr dla danego zestawu generatorów (\textbf{H2, H3, H5, H6}), dochodzimy w (\textbf{H9}) do następującej liczby odmiennych sposobów zamknięcia rezonansowych supersymetrycznych algebr:
\begin{table}[h!]
	\label{table-superalgebras}
	\begin{center}
		\begin{tabular}{l|l|l}
			~1$\times$ J             & ~1$\times$ J+Q             & ~~0$\times$ J+QY             \\ \hline
			~2$\times$ JP            & ~2$\times$ JP+Q            & ~\textbf{10}$\times$ JP+QY    \\ \hline
			~6$\times$ JPZ    & ~9$\times$ JPZ+Q  & \textbf{102}$\times$ JPZ+QY  \\ \hline
			30$\times$ JPZU & \textbf{43}$\times$ JPZU+Q &\textbf{667}$\times$ JPZU+QY
		\end{tabular}
		\caption[]{Rezonansowe algebry i superalgebry w zależności od generatorów.}
	\end{center}
\end{table}

Wszystkie uzyskane wyniki (super)algebr w (\textbf{H9}) są przedstawione w tabelach komutacyjnych, które stanowią zwięzły format natychmiast podkreślający wszystkie różnice między algebrami. Taka forma podkreśla również niezależność stałych struktury, dzięki czemu w zależności od konkretnego zastosowania możemy używać notacji stałych struktury AdS z przypadku $4D$ lub $3D$.

Aby zapewnić konstrukcję modeli supergrawitacji, wymagamy $\{Q, Q\}= P+...$ (z dodatkowymi wkładami lub bez nich). To ograniczenie znacząco redukuje uzyskane przypadki. Takie przykłady będziemy nazywać rezonansowymi algebrami supergrawitacji:
\begin{table}[h!]
	\label{table-supergravityalgebras}
	\begin{center}
		\begin{tabular}{l|l}
			~2$\times$ JP+Q & ~~~8$\times$ JP+QY    \\ \hline
			~6$\times$ JPZ+Q  & ~~71$\times$ JPZ+QY  \\ \hline
			17$\times$ JPZU+Q & 264$\times$ JPZU+QY
		\end{tabular}
		\caption[]{Możliwe rezonansowe superalgebry do konstruowania supergrawitacji.}
	\end{center}
\end{table}
\\
Poniżej znajduje się przegląd przykładów, które stoją za podanymi konfiguracjami:
\begin{itemize}
	\item Rezonansowe J, JP, JPZ, JPZU algebry
\end{itemize}
Konfiguracja $J$ odpowiada algebrze Lorentza, podczas gdy dwa przykłady $JP$ reprezentują algebry Poincar\'ego i AdS. Dodanie superładunku ($Q$) prowadzi do superalgebr $JP+Q$ Poincar\'ego i AdS. Przypadki $JPZ$ i $JPZU$, opisane w (\textbf{H3, H4, H5}), obejmują dobrze znane przykłady bozonowe, takie jak algebra Maxwella $\mathfrak{B}_{4}$ \cite{Bacry:1970ye, Schrader:1972zd} oraz $\mathfrak{C}_{4}$ wprowadzone przez Sorokę-Sorokę~\cite{Soroka:2006aj, Durka:2011nf, Salgado:2013eut}, a także ich uogólnienia do rodzin $\mathfrak{B}_{m}$ i $\mathfrak{C}_{m}$.
\begin{itemize}
	\item Rezonansowe J+Q, JP+Q, JPZ+Q superalgebry
\end{itemize}
Konfiguracja $J+Q$ odpowiada superalgebrze Lorentza. Przypadek $JP+Q$ obejmuje $\mathcal{N}=1$ supergrawitacyjne rozszerzenia algebr Poincar\'e i AdS. Niedawna praca (\textbf{H6}) zidentyfikowała konfiguracje $JPZ+Q$, jako rozszerzenia rezonansowe z dodatkowym generatorem $Z_{ab}$.
\begin{itemize}
	\item Rezonansowe JPZU+QY superalgebry
\end{itemize}
 Rozszerzenie zestawu generatorów wykonane w (\textbf{H9}) odtwarza znane przykłady i ujawnia nowe (super)algebry, wyróżnione pogrubieniem. Łącznie 667 superalgebr z podwojonymi generatorami podobnymi do Lorentza, translacji i superładunku ($2 \times (JP+Q)$), wynikających bezpośrednio z 30 bozonowych rodzin $JPZU$, wymagało ponad $10^{10}$ symbolicznych sprawdzeń (osiągniętych w ciągu jednego dnia dzięki niezwykle zoptymalizowanym obliczeniom opisanym w (\textbf{H9})). Spośród nich tylko 264 przypadki umożliwiają supergrawitację z $\{Q, Q\} = P+...$. Co istotne, jedynie trzy przypadki odtwarzają supersymetryczną podalgebrę AdS z $[P, P] \sim J$ i $\{Q, Q\} \sim P+J$, co ma potencjał dla przyszłych zastosowań.

%%%%%%%%%%%%%%%%%%%%%%%%%%%%%%%%%%%%%%%%%%%%%%%%%%%%%%%%%%%%%%%%%%%%%%%%%%%%%%%%%
\subsubsection{Rezonansowe Modele 3D i 4D}\label{Models}

Dla każdej z (super)algebr można skonstruować odrębny model (super)grawitacji zdefiniowany przez odpowiedni Lagranżjan/działanie z nowymi polami związanymi z nowymi generatorami. Jeśli $\{Q, Q\} \sim P+\ldots$ pojawia się w zamknięciu superalgebry, nazywamy je modelami supergrawitacyjnymi (\textbf{H6, H9}).

Niezbędne elementy tej konstrukcji to: (i) jednoforma cechowania~$\mathbb{A}$; (ii) dwuforma superkrzywizny~$\mathbb{F}$; (iii) parametr cechowania~$\Theta$, wraz z transformacjami cechowania; (iv) forma Killinga $\left\langle \dots \right\rangle$ zapewniająca ślad po indeksach grupy.

\begin{itemize}
	\item Dwuforma superkrzywizny jest bezpośrednio zbudowana z pola cechowania $\mathbb{A}$, jako $\mathbb{F}=d\mathbb{A}+\frac{1}{2}\left[[\mathbb{A},\mathbb{A}\right]]$, gdzie
	używamy określonej algebry do obliczenia (anty)komutatorów.
	\item Rozszerzony parametr cechowania, na przykład $\Theta=\frac{1}{2} \lambda^{ab}J_{ab}+ \tau^{a}P_{a}+\frac{1}{2} \tilde{\lambda}^{ab}Z_{ab}+\tilde{\tau}^{a} U_{a}+\epsilon^{\alpha}Q_{\alpha}+\varepsilon^{\alpha}Y_{\alpha}$,
	pozwala nam zapisać prawo transformacji dla konkretnych pól względem Lorentza $\lambda^{ab}$, translacji $\tau^a$, tzw. transformacji „Maxwellowskich”: $\tilde{\lambda}^{ab}$, $\tilde{\tau}^a$ oraz obu superładunków $\epsilon^\alpha$, $\varepsilon^\alpha$.
	\item Prawo transformacji $\delta_\Theta \mathbb{A}= \mathbb{D}^\mathbb{A}\Theta= d\Theta+[[\mathbb{A},\Theta]]$ opisuje szczególne prawa jednoznacznie określone przez konkretną (super)algebrę, z odpowiednimi:
	$$\delta_\Theta\omega^{ab},\delta_\Theta e^{a},\delta_\Theta h^{ab},\delta_\Theta k^{a},\delta_\Theta\psi,\delta_\Theta\chi$$
	transformacjami. 
	\item Dodatkowo zapisujemy również transformację krzywizny $\delta_\Theta \mathbb{F}= [[\Theta,\mathbb{F}]]$.
\end{itemize}
W zależności od konstrukcji możemy osiągnąć:
\begin{itemize}
	\item w $3D$ dla teorii Cherna-Simonsa \cite{Concha:2014tca, Hassaine:2016amq}) pełną niezmienniczość działania względem symetrii cechowania,
	\item w $4D$ ograniczenie tylko do niezmienniczości względem Lorentza i supersymetrii, z możliwością dodatkowej niezmienniczości wynikającej z „symetrii Maxwellowskich” \cite{Durka:2011gm, Durka:2011va, Kamimura:2011mq}.
\end{itemize}
W przypadku $4D$, aby zapewnić supersymetrię, należy zastosować formalizm $1.5$ rzędu. W związku z tym konstrukcja w $4D$ jest bardziej złożona (zobacz metody już zastosowane w zdeformowanej teorii BF \cite{Durka:2011gm, Durka:2012wd} oraz modelach MacDowella-Mansouriego \cite{Gocanin:2019ota, MacDowell:1977jt}). To sprawia, że działanie Cherna-Simonsa (CS) jest znacznie częściej stosowane w badaniach.

\subsubsection{Model Cherna-Simonsa i podsektory niezmiennicze}

W $3D$ działanie Cherna-Simonsa (CS)
\begin{align}\label{CS}
	I_{CS} = \frac{k}{4\pi} \int_{\mathcal{M}} \left\langle \mathbb{A} \wedge d\mathbb{A} + \frac{1}{3} \mathbb{A} \wedge [\mathbb{A}, \mathbb{A}] \right\rangle
\end{align}
zapewnia wygodne ramy z prostą konstrukcją działania, gwarantując pełną niezmienniczość symetrii. Dla niezerowej stałej kosmologicznej \(\Lambda = \mp \frac{(D-1)(D-2)}{2\ell^{2}}\) oraz pola cechowania (A)dS $A = \frac{1}{2}\omega^{ab} J_{ab} + \frac{1}{\ell} e^{a} P_{a}$, działanie grawitacyjne w 3D przyjmuje postać:
\begin{align}
	I^{3D}_{e,\omega}&=\frac{1}{16\pi G}\int \left[ \alpha _{0}\left(
	\omega^{ab}\wedge d\omega _{ab}+\frac{1}{3}\omega^{ab} \wedge \omega
	_{bc}\wedge \omega^{c}{}_{a}+\frac{1}{\ell^2}e_{a} \wedge D_{\omega }e^{a} \right) \right.  \notag \\
	&\left. +\alpha_{1}\left(\mathcal{R}^{ab}(\omega)\wedge e^{c}+\frac{1}{3\ell^2}e^{a} \wedge e^{b} \wedge e^{c}\right) \epsilon
	_{abc}\right] \,.
\end{align}

Należy zauważyć, że w przypadku trójwymiarowym z $a,b=0,1,2$ można wprowadzić dualne generatory, takie jak $J_{a}=\frac{1}{2}\epsilon_{a}{}^{bc}\,J_{bc}$, co pozwala na jednolite zapisanie AdS:
\begin{align}
	\left[ J_{a},J_{b}\right] =\epsilon_{abc}J^{c},\quad
	\left[ J_{a},P_{b}\right] =\epsilon _{abc}P^{c},\quad
	\left[ P_{a},P_{b}\right] =\epsilon _{abc}J^{c}\,.
\end{align}
z $\epsilon$ jako antysymetrycznym symbolem Levi-Civity.

%%%%%%%%%%%%%%%%%%%%%%%%%%%%%%%%%%%%%%%%%%%%%%%%%%%%%%%%%%%%%%%%%%%%%%%%%%
Jedną z najważniejszych cech wykazywanych przez struktury rezonansowe/półgrupowe (\textbf{H3, H4, H6}) jest istnienie podsektorów niezmienniczych poprzez dowolnie wartościowane stałe $\alpha$ mnożące tensory niezmiennicze.

Metryka Killinga, używana do kontrakcji wszystkich indeksów grupowych, przyjmuje formę tensora niezmienniczego określonego dla dowolnej kombinacji dwóch generatorów. Na przykład dla superalgebry $\mathfrak{C}_4$ Soroki-Soroki zapisujemy $\left\langle J_{ab}\,J_{cd}\right\rangle =\alpha_{0}\,\epsilon_{abcd}$, $\left\langle P_{a}\,P_{b}\right\rangle =\alpha_2 \eta_{ab}$, $\left\langle J_{ab}\,Z_{cd}\right\rangle =\alpha_{2}\epsilon_{abcd}$, $\left\langle Q_\alpha\, Q_\beta\right\rangle=(\alpha_1-\alpha_2)\,C_{\alpha\beta}$ i tak dalej. Składowe różne od zera dla danej (super)algebry mogą być ustalone na podstawie znanych wyników super-AdS oraz tożsamości $\langle[[X_i, X_j ]]\, X_k \rangle =\langle X_i \, [[X_j, X_k]] \rangle$, gdzie $X_i$ jest dowolnym generatorem. Mogą one również pochodzić z ekspansji półgrupy. Ich pojawienie się jest bezpośrednio związane z formą konkretnych (super)algebr (\textbf{H3, H4, H5, H6, H9}), co prowadzi do różnorodnej zawartości i podziału członów między różne podsektory.

Jako przykład rozważmy kilka konfiguracji $JPZ+Q$ (\textbf{H6}) z przejściem do dualnych pól $3D$ $\omega^a=\frac{1}{2}\epsilon^{abc}\omega_{bc}$, $h^a=\frac{1}{2}\epsilon^{abc}h_{bc}$ oraz definicji generatorów, $J_a=\frac{1}{2}\epsilon_{a}{}^{bc}J_{ab}$, $Z_a=\frac{1}{2}\epsilon_{a}^{bc}{}Z_{ab}$. Wprowadźmy również $\mathcal{R}^{a}=d\omega^{a}+\frac{1}{2}\epsilon^{abc}\omega_{b}\omega_{c}$, $D_{\omega}e^{a}=d e^{a}+\epsilon^{abc}\omega_{b}\,e_{c}$, a także kowariantną pochodną Lorentza działającą na spinor $\mathcal{D}_{\omega }\psi=d\psi+\frac{1}{2}\omega^{a}\gamma_{a}\psi$.

Najbardziej ogólne trójwymiarowe działanie CS supergrawitacji $\mathcal{N}=1$, Eq.~\eqref{CS}, które jest niezmiennicze względem superalgebry Soroki-Soroki $\mathfrak{C}_{4}$, przyjmuje postać:
\begin{eqnarray}
	I_{CS}^{\mathfrak{C}_{4}} &=&\frac{k}{4\pi }\int \left[ \alpha _{0}\left(
	\omega ^{a}d\omega _{a}+\frac{1}{3}\epsilon ^{abc}\omega _{a}\omega
	_{b}\omega _{c}\right) \right.  \notag \\
	&&\left. +\alpha _{1}\left( 2\mathcal{R}^{a}e_{a}+\frac{1}{3\ell^2}\epsilon
	^{abc}e_{a}e_{b}e_{c}+2 e_{a}D_{\omega }h^{a}+\epsilon ^{abc}e_{a}h_{b}h_{c}+2
	\bar{\psi}\mathcal{F} \right) \right.  \label{SSCS} \\
	&&+\alpha _{2}\left(\frac{1}{\ell^2}e_{a}D_{\omega }e^{a}+2h_{a}\mathcal{R}
	^{a}+\frac{1}{\ell^2}\epsilon^{abc}e_{a}e_{b}h_{c}\notag\left. +h_{a}D_{\omega }h^{a}+\frac{1}{3}\epsilon ^{abc}h_{a}h_{b}h_{c}-\frac{2}{\ell}\bar{\psi}\mathcal{F} \right) \right] \,,  \label{CSmathfrak_C4}
\end{eqnarray}
where $\mathcal{F}=\mathcal{D}_{\omega }\psi+\frac{1}{2\ell}e^{a}\gamma_{a}\psi+\frac{1}{2}h^{a}\gamma_{a}\psi$. 

Jednocześnie dla algebry Maxwella $\mathfrak{B}_4$ działanie CS wynosi:
\begin{eqnarray}\label{CSmathfrak_B4} 
	I_{CS}^{\mathfrak{B}_{4}} &=&\frac{k}{4\pi }\int \left[ \alpha _{0}\left(
	\omega ^{a}d\omega _{a}+\frac{1}{3}\epsilon ^{abc}\omega _{a}\omega
	_{b}\omega _{c}\right)\right. \notag\\
	&& +\alpha _{1}\left( 2\mathcal{R}^{a}e_{a}\right)\notag\\
	&&\left. \alpha _{2}\left(\frac{1}{\ell^2}
	e_{a}D_{\omega }e^{a}+2h_{a}\mathcal{R}^{a}-\frac{2}{\ell}\bar{\psi}\mathcal{D}_{\omega }\psi\right) \right] \,,
\end{eqnarray}
podczas gdy przypadek $B_4$ prowadzi do działania postaci:
\begin{eqnarray}
	I_{CS}^{B_{4}} &=&\frac{k}{4\pi }\int \left[ \alpha _{0}\left(
	\omega ^{a}d\omega _{a}+\frac{1}{3}\epsilon ^{abc}\omega _{a}\omega
	_{b}\omega _{c}+h_{a}D_{\omega}h^{a}\right)   \right.  \notag \\
	&&+\alpha _{1}\left(2\mathcal{R}^{a}e_{a}+2e_{a}D_{\omega }h^{a}+\epsilon ^{abc}e_{a}h_{b}h_{c}+2
	\bar{\psi}\mathcal{D}_\omega \psi+\frac{1}{\ell} \bar{\psi}h_{a}\gamma^{a} \psi \right)\notag\\
	&&\left.  +\alpha _{2}\left(2h_{a}\mathcal{R}
	^{a}+\frac{1}{3}\epsilon ^{abc}h_{a}h_{b}h_{c} \right) \right]\,.  \label{CS_B4} 
\end{eqnarray}
Standardowy przypadek działania CS niezmienniczego względem superalgebry AdS w trójwymiarowej supergrawitacji $\mathcal{N}=1$ ma postać:
\begin{eqnarray}
	I_{CS}^{AdS} &=&\frac{k}{4\pi }\int \left[ \alpha _{0}\left(
	\omega ^{a}d\omega _{a}+\frac{1}{3}\epsilon^{abc}\omega _{a}\omega
	_{b}\omega _{c}+\frac{1}{\ell^2}e_{a}D_{\omega }e^{a}-\frac{2}{\ell}\bar{\psi}\mathcal{D}_{\omega }\psi-\frac{1}{\ell^2}\bar{\psi}e^{a}\gamma_{a}\psi\right) \right.  \notag \\
	&& +\alpha _{1}\left( 2 \mathcal{R}^{a}e_{a}+\frac{1}{3\ell^2}\epsilon
	^{abc}e_{a}e_{b}e_{c}+2
	\bar{\psi}\mathcal{D}_{\omega }\psi+\frac{1}{\ell}
	\bar{\psi}e^{a}\gamma_{a}\psi \right) \label{CSAdS}\,.
\end{eqnarray}
Po rozważeniu granicy $\ell \rightarrow \infty $, wynikowe działanie CS opisuje trójwymiarową supergrawitację $\mathcal{N}=1$ Poincar\'e:
\begin{eqnarray}
	I_{CS}^{Poincare}&=&\frac{1}{16\pi G }\int \left[ \alpha_{0}\left( \omega
	^{a}d\omega _{a}+\frac{1}{3}\epsilon ^{abc}\omega _{a}\omega _{b}\omega
	_{c} \right)\right.  \notag \\
	& &\left.+\alpha _{1}\left( 2\mathcal{R}^{a}e_{a}+2\bar{\psi} \mathcal{D}_{\omega}\psi \right) \right] \,.
\end{eqnarray}

Ponadto prowadzi to do istotnych konsekwencji omówionych w (\textbf{H6}): w ramach podsektorów niezmienniczych obserwujemy albo reorganizację członów, albo ich brak. To z kolei naturalnie prowadzi do wyraźnie różnych (nadmiernie ograniczających) równań pola.

Należy zauważyć, że działanie CS supergrawitacji zawiera trzy niezależne sektory proporcjonalne do $\alpha _{0}$, $\alpha _{1}$ i $\alpha _{2}$. Człon proporcjonalny do $\alpha _{1}$ zawsze zawiera człon Einsteina-Hilberta. Nowy wkład jest proporcjonalny do $\alpha _{2}$. Można wybrać pewne stałe $\alpha$ równe zeru i skupić się wyłącznie na podsektorach niezmienniczych. Na przykład w przypadku (\textbf{H4}) podejście to polegało na rozważaniu bozonowych działań z $\alpha_1 = 0$.

%%%%%%%%%%%%%%%%%%%%%%%%%%%%%%%%%%%%%%%%%%%%%%%%%%%%%%%%%%%%%%%%%%%%%%%%%%
\subsubsection{Model MacDowella-Mansouriego}

Model MacDowella-Mansouriego formułuje grawitację w czterech wymiarach, osadzając jej stopnie swobody w większej strukturze cechowania, mianowicie symetrii de Sittera (dS) lub Anty-de Sittera (AdS) \cite{Durka:2012wd, MacDowell:1977jt}. Następnie ta symetria cechowania jest łamana w celu uzyskania Ogólnej Teorii Względności.

Pełna koneksja AdS $so(2,3)$-wartościowa $\mathbb{A}^{IJ}$ (wybrana ze względu na późniejsze zastosowania w supergrawitacji) łączy vierbein $e^a$ i koneksję spinową $\omega^{ab}$,
\begin{equation}
	A^{ab}=\omega^{ab}\,,\qquad  A^{a4}=\frac{1}{\ell}e^a\,,
\end{equation}
gdzie $I,J=0,1,2,3,4$ są indeksami AdS, a $a, b = 0, 1, 2, 3$ są indeksami Lorentza. Działanie grawitacyjne jest zbudowane przy użyciu dwuformy krzywizny analogicznie do teorii Yang-Millsa, z dodatkowym wprowadzeniem mechanizmu łamania symetrii, który redukuje pierwotną symetrię cechowania do grupy Lorentza. Nie możemy użyć pełnej krzywizny $F^{IJ}(A)$ do budowy działania teorii dynamicznej w czterech wymiarach. Takie sformułowanie jest jednak możliwe przy użyciu grupowej dualności ($\star$) i złamaniu symetrii cechowania przez projekcję pełnej krzywizny na indeksy Lorentza.
\begin{equation}
	\mathbb{F}^{IJ}\quad\rightarrow\quad \hat{\mathbb{F}}^{IJ}= F^{ab}\quad\quad \mathrm{where}\quad\quad F^{ab}=\mathcal{R}^{ab}+\frac{1}{\ell^2}e^a\wedge e^b\,.
\end{equation} 
Ogólna Teoria Względności postrzegana jako teoria łamania symetrii cechowania wyłania się wówczas z działania
\begin{equation}
	I_{MM}(A)=\frac{\ell^2}{64\pi G} \int tr\big(\hat{\mathbb{F}} \wedge \star \hat{\mathbb{F}}\big)\,,
\end{equation}
\begin{equation}
	I_{MM}(A)=\frac{\ell^2}{64\pi G} \int \left(\mathcal{R}^{ab}+\frac{1}{\ell^2}e^a\wedge e^b\right)\wedge \left(\mathcal{R}^{cd}+\frac{1}{\ell^2}e^c\wedge e^d\right)\epsilon_{abcd}\,.
\end{equation}
Ta intrygująca teoria grawitacji, oparta na grupie (A)dS, pod spodem sprowadza się do
\begin{align}
	32\pi G\,I_{MM} &= \int \epsilon_{abcd}\mathcal{R}^{ab}\wedge e^{c}\wedge \,e^{d}+\frac{1}{\,2\ell^2}\int\epsilon_{abcd} e^{a}\,\wedge e^{b}\wedge e^{c}\,\wedge
	e^{d} +\frac{\ell^2}{2}\int\epsilon_{abcd}\mathcal{R}^{ab}\wedge \mathcal{R}^{cd}\,.\notag\\
	32\pi G\, I_{MM} &=\int \mathrm{Einstein/Cartan}+\frac{1}{\,2\ell^2}\int  \mathrm{cosmological}+\frac{\ell^2}{2}\int \mathrm{Euler} \,, 
\end{align}
i wiąże stałą kosmologiczną z niezmiennikiem Eulera (członem kwadratowym w formach krzywizny). Równania pola pozostają takie same jak w pierwszorządowym działaniu Einsteina-Cartana Eq.~\eqref{Einstein}, ponieważ wariacja względem członu Eulera prowadzi do $D^\omega \mathcal{R}^{ab}$, które znika z powodu tożsamości Bianchiego.

Co więcej, działanie MacDowella-Mansouriego w naturalny sposób uwzględnia supersymetryczne rozszerzenie. Wersja SUGRA pojawia się w sposób naturalny, ponieważ do bozonowej części (w wersji supersymetrycznej) dodawany jest fermionowy odpowiednik $\sim F^{(s)ab} \epsilon_{abcd} F^{(s)cd} + \bar{\mathcal{F}}\wedge\gamma^5\mathcal{F}$, gdzie krzywizna fermionowa $\mathcal{F} = \mathcal{D}^A\psi=\mathcal{D}^\omega\psi+\frac{1}{2\ell}
e^{a}\,\gamma_{a}\,\psi= 
d\psi+\frac{1}{4}\omega^{ab}\,\gamma_{ab}\,\psi+\frac{1}{2\ell}
e^{a}\,\gamma_{a}\,\psi$.

%%%%%%%%%%%%%%%%%%%%%%%%%%%%%%%%%%%%%%%%%%%%%%%%%%%%%%%%%%%%%%%%%%%%%%%%%%
\subsubsection{Teoria BF MacDowella-Mansouriego}

W przypadku teorii BF typu MacDowella-Mansouriego \cite{Freidel:2005ak} (często nazywanej również zdeformowaną), punktem wyjścia jest topologiczne działanie $BF$ z dodatkowym członem deformacyjnym, który łamie symetrię cechowania do lokalnej symetrii Lorentza, analogicznie do schematu MacDowella-Mansouriego:
\begin{align}
	32\pi G\, I_{BF} (B,A)&= \int tr(\mathbb{B}\wedge \mathbb{F}-\frac{\beta}2\, \mathbb{B}\wedge \mathbb{B}- \frac\alpha4\,\, \hat{\mathbb{B}}\wedge \hat{\mathbb{B}})\notag\\
	&= \int B_{IJ}\wedge F^{IJ}(A) -\frac{\beta}2\, B_{IJ}\wedge B^{IJ} - \frac\alpha4\, \epsilon_{abcd}\, B^{ab}\wedge B^{cd}\,,
\end{align}
gdzie $F^{IJ}(A)$ jest krzywizną koneksji AdS $A^{IJ}$, a indeksy $I,J, \ldots$ przyjmują wartości od 0 do 4, $I=(a,4)$. Stanowi to kontrast względem konstrukcji Plebańskiego \cite{Plebanski:1977zz}, gdzie uwzględnia się tylko wiąz $B = (\gamma^{-1} +\star)\, e\wedge e$.

Dwuforma pomocniczego pola $B$ po rozwiązaniu równań pola może być rozłożona w następujący sposób:
\begin{align}
	B^{ab}= \frac{\alpha}{(\alpha^{2}+\beta^{2})}\left(\frac{1}{2}\frac{\beta}{\alpha}\,\delta^{ab}_{cd}- \frac{1}{2}\epsilon^{ab}{}_{cd}\right)F^{cd}\,,\qquad 	
	B^{a4} =  \frac{1}{\beta}\,F^{a4}\,,
\end{align}
gdzie składniki krzywizny $F^{IJ}=d A^{IJ}+A^{I}{}_{K}\wedge A^{KJ}$ to krzywizna Lorentza $\mathcal{R}^{ab}(\omega)$ z uwzględnieniem stałej kosmologicznej oraz torsja $T^a=D^\omega e^a$:
\begin{align}
	F^{ab}=\mathcal{R}^{ab}(\omega )+\frac{1}{\ell ^{2}}\,e^{i}\wedge e^{j}\,,\qquad	
	F^{a4}=\frac{1}{\ell }T^{a}\,.
\end{align}
Podstawiając te równania do pierwotnego Lagranżjanu z $I_{BF}=\int L_{BF}$ i używając jawnych wyrażeń dla parametrów:
\begin{align}\label{constant}
	\alpha= \frac{ \Lambda}{3}\frac{\gamma^2}{1+{\gamma^2}}\,,\qquad {\beta}=-\frac{\alpha}\gamma\,,\qquad \Lambda=- \frac{3}{\ell^2} \,,
\end{align}
widzimy, że ostatecznie teoria BF MacDowella-Mansouriego \cite{Durka:2012wd, Freidel:2005ak} skrywa najbardziej ogólny Lagrangian pierwszego rzędu:
\begin{align*}\label{full_action}
	32\pi G\,\mathcal{L}_{BF}& = \epsilon_{abcd}\mathcal{R}^{ab}\wedge e^{c}\wedge \,e^{d}-\frac{\Lambda}{6}\epsilon_{abcd} e^{a}\,\wedge e^{b}\wedge e^{c}\,\wedge
	e^{d}  -{2}{\gamma }\mathcal{R}_{ab}\wedge e^{a}\wedge e^{b} \nonumber \\
	& -2 \frac{1+\gamma^2}{\gamma }\,\left(T^{a}\wedge T_{a}-\mathcal{R}_{ab}\wedge e^{a}\wedge
	e^{b}\right) -\frac{3}{\Lambda\gamma}\mathcal{R}^{ab}\wedge \mathcal{R}_{ab}-\frac{3}{2\Lambda}\epsilon_{abcd}\mathcal{R}^{ab}\wedge \mathcal{R}^{cd}\,.
\end{align*}
Pierwszy wiersz (razem z członem Eulera) odtwarza Lagranżjan MacDowella-Mansouriego $\mathcal{L}_{MM}$ oraz Lagranżjan Holsta \cite{Holst:1995pc}, gdzie stała sprzężenia jest odwrotnością parametru Barbero-Immirziego. Parametr Barbero-Immirziego odgrywa kluczową rolę w Pętlowej Grawitacji Kwantowej (LQG), służąc jako współczynnik członu Holsta \cite{Holst:1995pc} i pojawiając się w procesie kwantyzacji. Druga linia to kombinacja członów topologicznych: form Nieh-Yana, Pontryagina oraz Eulera. Można połączyć drugi człon niezmiennika Nieh-Yana z członem Holsta w pierwszej linii, aby ostatecznie uzyskać poprawny współczynnik $2/\gamma$. Warto zauważyć, że rozważany tutaj Lagranżjan może być wygodnie przepisany w stylu MacDowella-Mansouriego/Yang-Millsa jako suma kwadratów torsji i krzywizny AdS:
\begin{align}
	32\pi G\,\mathcal{L}_{BF}
	&=-\frac{3}{2\Lambda}\,F^{ab}\wedge\left(\frac{1}{\gamma}\delta_{abcd}+\epsilon_{abcd}\right)F^{cd}-2 \frac{1+\gamma^2}{\gamma }\,T^{a}\wedge T_{a}\,.
\end{align}
Taka ogólność, szczegółowo zbadana w \cite{Durka:2012wd}, jest głównym powodem jej stosowania w różnych aplikacjach (\textbf{H7, H8, H11}).

%%%%%%%%%%%%%%%%%%%%%%%%%%%%%%%%%%%%%%%%%%%%%%%%%%%%%%%%%%%%%%%%%%%%%%%%%%
\subsection{Mój wkład w badania}

%%%%%%%%%%%%%%%%%%%%%%%%%%%%%%%%%%%%%%%%%%%%%%%%%%%%%%%%%%%%%%%%%%%%%%%%%%%
\subsubsection{Rezonansowe (super)algebry i grawitacja (H2-6, H9, H10, H12)}

Pierwsza część moich badań koncentrowała się na wprowadzeniu, klasyfikacji i zastosowaniach rezonansowych (super)algebr — nowego podejścia do powiększania struktur algebraicznych, wykraczającego poza tradycyjne (super)algebry Poincarégo i AdS. Algebry te rozszerzają znane struktury poprzez włączenie dodatkowych generatorów, wzbogacając ramy algebraiczne i oferując nowe możliwości eksploracji teoretycznej. 

\myparagraph{New Family of Maxwell-like Algebras (H2)}

To otwierające badanie jest wynikiem mojej współpracy w ramach projektu podoktorskiego 2013-2016, nawiązanego z badaczami w Chile, którzy również interesowali się tematem tzw. algebry Maxwella \cite{Bacry:1970ye, Schrader:1972zd}. Moim głównym wkładem było wprowadzenie alternatywnej metody zamykania algebr podobnych do Maxwella, co doprowadziło do nowej rodziny algebraicznej oznaczonej jako $\mathfrak{D}$, będącej sumą prostą algebr Anti-de Sittera i Maxwella. Podejście to, dobrze osadzone w ramach S-ekspansji \cite{Izaurieta:2006zz}, pozwoliło na bezpośrednią konstrukcję teorii grawitacji opartych na nowo zidentyfikowanych rozszerzeniach.

Nowa rodzina została wprowadzona poprzez dodanie nowych generatorów (podobnych do Lorentza i translacji) $\{ Z_{ab}, R_{a}, \ldots \}$ do generatorów Lorentza i translacji. Ten nowy schemat różnił się od wcześniej znanych rodzin Maxwella: ($\mathfrak{B}_m$, tzw. Maxwellowskiej \cite{Bacry:1970ye, Schrader:1972zd}, oraz $\mathfrak{C}_m$, tzw. Soroki-Soroki \cite{Soroka:2006aj, Soroka:2011tc, Salgado:2013eut}, lub AdS-Maxwell \cite{Durka:2011nf}). Co ciekawe, $\mathfrak{D}_m$ w ramach kontrakcji Inönü-Wignera prowadzi do tradycyjnych algebr Maxwella $\mathfrak{B}_m$.

Algebry $\mathfrak{D}_m$ znalazły natychmiastowe zastosowanie w teoriach grawitacyjnych, dzięki bezpośredniej konstrukcji działań Cherna-Simonsa \eqref{CS} oraz Born-Infelda \cite{Concha:2014vka} w różnych wymiarach czasoprzestrzeni. Te teorie są budowane z użyciem odpowiednio rozszerzonej koneksji, co prowadzi do Lagranżjanów, które prosto łączą sektor grawitacyjny z dodatkowymi polami związanymi z nowymi generatorami. Takie zintegrowane podejście wzbogaca zawartość pól w modelach grawitacyjnych, umożliwiając nowe teoretyczne spostrzeżenia i praktyczne konstrukcje. Chociaż suma prosta może wydawać się nieinteresująca, należy pamiętać, że konfiguracja wciąż pozostaje otwarta na pewne ansatze lub łamanie symetrii.

%%%%%%%%%%%%%%%%%%%%%%%%%%%%%%%%%%%%%%%%%%%%%%%%%%%%%%%%%%%%%%%%%%%%%%%%%%%
\myparagraph{Resonant algebras and gravity (H3)} 

Kolejna, niezależna, praca opiera się na istniejącym zrozumieniu rozszerzeń algebraicznych w algebrach Liego. W szczególności koncentrując się na algebrze Maxwella i wykracza poza nią, bezpośrednio przedstawia moją próbę zmiany tradycyjnego podejścia do S-rozszerzonych algebr Liego, szczególnie po wynikach (\textbf{H2}), które ujawniły wcześniej nieodkryte rozszerzenia. Ta praca stanowi istotny postęp w tej dziedzinie, przesuwając ramy S-ekspansji w kierunku eksploracji maksymalnej swobody w zamykaniu tabel mnożenia dla abelowych półgrup.

Kluczową innowacją w tych badaniach było uwzględnienie elementu zerowego w dekompozycji rezonansowej, wraz z powiązaniem generatora Lorentza z elementem neutralnym półgrupy. Taka zmiana metodologii prowadzi do szerokiej klasy rozszerzonych algebr Liego, jednocześnie wykazując użyteczność w generowaniu nowych teorii grawitacji i supergrawitacji.

W tej pracy wprowadzono uogólnienie koncepcji \textit{rezonansowych algebr}, zapewniając kompleksową klasyfikację do pewnej liczby generatorów. Praca ukazała potencjał tych algebr w dostarczaniu alternatywnych sformułowań grawitacyjnych teorii cechowania, oferując szerszy zakres struktur algebraicznych. Badanie z powodzeniem odtworzyło wszystkie algebry Maxwella typu $\mathfrak{B}_m$ \cite{Schrader:1972zd}, $\mathfrak{C}_m$ \cite{Soroka:2006aj, Salgado:2013eut} oraz wcześniej wprowadzoną $\mathfrak{D}_m$ (\textbf{H2}), jak również zaproponowało nowe przykłady. Wyczerpano tym samym generalizację bozonowych rozszerzeń do określonej liczby generatorów podobnych do Lorentza i translacji.

Dla zbioru generatorów $J_{ab}, P_{a}, Z_{ab}$, wyniki pracy obejmowały cztery nowe algebry podobne do Poincarégo, oprócz dwóch znanych algebr podobnych do Maxwella ($\mathfrak{B}_4$ i $\mathfrak{C}_4$). Dodatkowo okazało się, że jedna ze znalezionych algebr podobnych do Poincarégo mogła być zapisana jako suma prosta algebr AdS i Lorentza, co otwiera nowe kierunki badań w tematach, gdzie rozważano sumę prostą AdS i Poincarégo.

Oprócz tych osiągnięć, w pracy (\textbf{H3}) podkreśliłem bogactwo pomijanych gałęzi rozszerzeń, gdy do generatorów $\{J, P, Z\}$ dodany zostanie nowy generator translacji $R_{a}$. Praca wskazuje istotne uogólnienie dotyczące porządku rozszerzeń, takie jak włączenie $R_a$ do $J_{ab}, P_a$ zamiast $Z_{ab}$. Zwraca również uwagę na szczególną rolę $J_{ab}$, który nie może być zamieniony z $Z_{ab}$, ponieważ zawsze chcemy zachować niezmienniczość względem Lorentza.

Poza strukturą algebraiczną, praca ta dostarcza wszystkich niezbędnych składników do konstrukcji działań grawitacyjnych opartych na uzyskanych wynikach. W przypadkach $(2n+1)$-wymiarowych można zastosować działanie grawitacji Cherna-Simonsa, natomiast w scenariuszach $2n$-wymiarowych można wykorzystać działanie typu Born-Infelda \cite{Concha:2014vka} lub eksplorować formułowania w $4D$ za pomocą teorii BF \cite{Durka:2012wd, Durka:2011gm}.

%%%%%%%%%%%%%%%%%%%%%%%%%%%%%%%%%%%%%%%%%%%%%%%%%%%%%%%%%%%%%%%%%%%%%%%%%%%%

\myparagraph{Resonant algebras in CS model of topological insulators (H4)} 

Tutaj zastosowałem rezonansowe algebry z (\textbf{H3}) do badania izolatorów topologicznych, demonstrując ich użyteczność w konstrukcji działań/modeli. Podejście to, dalekie od zastosowań w modelach grawitacyjnych, wprowadza nowe $3D$ sformułowania Lagranżjanów zdolnych do opisu różnych efektów fizycznych związanych z zjawiskiem efektu Halla. Podążając za pracą G.~Palumbo \cite{Palumbo:2016nku}, możemy dostarczyć systematyczne mateamtyczne ramy do opisu izolatorów topologicznych oparte na sześciu rezonansowych algebrach związanych z generatorami $J_{ab}$, $P_a$ oraz $Z_{ab}$.

Pionierska praca Palumbo podkreśliła kluczową rolę modelu Cherna-Simonsa jako efektywnej teorii pola dla układów Kwantowego Efektu Halla oraz jego zdolność do uchwycenia unikalnych zjawisk fizycznych w izolatorach topologicznych. Te intrygujące materiały, mimo że są izolatorami w swojej objętości, wykazują nietrywialne właściwości przewodzące na swoich krawędziach.

Człon Wen-Zee \cite{Wen:1992ej}, sprzęgający potencjał elektromagnetyczny $A_\mu$ z koneksją spinową $O(2)$ $\widehat{\omega}$, wzbogaca teorię, ale jest wprowadzany heurystycznie, a nie wyprowadzany z zasad teorii CS. Wskazuje to na lukę w jednoczesnym uwzględnieniu działań "grawitacyjnych" ($\omega$) i elektromagnetycznych ($A$), ponieważ brakuje im naturalnej interakcji wykraczającej poza minimalne sprzężenie znane z grawitacji sprzężonej z elektromagnetyzmem. Algebry Maxwella typu $\mathfrak{B}$ rozwiązują ten problem, wprowadzając generator $Z_{ab}$ oraz powiązane pole $h^{ab}_\mu$, które integruje $A_\mu$ w ramach swojej definicji.

Konstrukcja działania zaproponowana przez Palumbo \cite{Palumbo:2016nku} zawierała jednak krytyczną wadę: zastosowana algebra Maxwella $\mathfrak{B}$ była niewystarczająca do wygenerowania wszystkich wymaganych członów dla spójnego sformułowania. W przeciwieństwie do tego, moje badanie (\textbf{H4}) koncentruje się głównie na algebrze $\mathfrak{C}_4$, zdefiniowanej w $3D$, ponieważ oferuje ona najbogatszą zawartość członów w (topologicznym) Lagranżjanie spośród sześciu rezonansowych algebr. Głównym celem było umożliwienie zunifikowanego opisu standardowych i potencjalnie niestandardowych izolatorów topologicznych poprzez naturalną konstrukcję działań opartą na rezonansowych algebrach $JPZ$. Obejmowało to włączenie relatywistycznego członu Wen-Zee, wraz z innymi członami, z wartościami stałych sprzężenia automatycznie ustalonymi przez algebrę.

Poprzez wyznaczenie wszystkich sześciu rezonansowych algebr utworzonych przez generatory Lorentza, translacji oraz Maxwella, wyprowadziliśmy sześć odrębnych Lagranżjanów Cherna-Simonsa, z których każdy prowadzi do zmodyfikowanych opisów izolatorów topologicznych. Badanie kończy się wnioskiem, że przykład $\mathfrak{C}_4$ jest najbardziej odpowiedni dla izolatorów topologicznych, ponieważ posiada nieznikającą torsję związaną z lepkością w efekcie Halla. Konfiguracje algebr takie jak $B_4,\tilde{B}_4$ i $\mathfrak{B}_4$ nie zawierają członu $AdA$ (a zatem brak w nich normalnych odpowiedzi prądowych '$E$' i '$B$' dla efektu Halla), podczas gdy, co interesujące, nadal obecny jest człon Wen-Zee. Niewątpliwie cenne będzie ponowne i dalsze zbadanie tych intrygujących efektów Halla.

Niezmiennicze metryki (Killinga) dla tych algebr zostały zidentyfikowane poprzez formy biliniowe z udziałem dowolnych rzeczywistych parametrów $\alpha_0$ i $\alpha_2$. To sformułowanie ułatwiło konstrukcję zwięzłych modeli Cherna-Simonsa z dobrze zdefiniowanymi stałymi i równaniami pola, oferując prostszą ścieżkę dla praktycznych zastosowań. Podejście to kontrastuje z metodami opisanymi w \cite{Palumbo:2016nku}, gdzie konstrukcja była przeprowadzana ręcznie, bez uzasadnienia mieszając cechy algebr $\mathfrak{B}$ i $\mathfrak{C}$.

Wraz z tym, w artykule (\textbf{H4}) zaproponowaliśmy również uzupełniające (grawitacyjne) działania oparte na alternatywnej niezmienniczej postaci metryki Killinga $\epsilon_{abc}$, która może być potencjalnie interesująca dla uogólnionych $3D$ modeli grawitacyjnych.

%%%%%%%%%%%%%%%%%%%%%%%%%%%%%%%%%%%%%%%%%%%%%%%%%%%%%%%%%%%%%%%%%%%%%%%%%%%%	
\myparagraph{On the Number of Possible Resonant Algebras (H5)} 

W ramach tego badania przeprowadziłem kolejne obliczeniowe studium, mające na celu sklasyfikowanie możliwych bozonowych rezonansowych algebr dla dalszych rozszerzonych konfiguracji generatorów, wykraczając poza te zaprezentowane w (\textbf{H3}). Powstała praca systematycznie zbadała konstrukcję i klasyfikację rezonansowych algebr (nadal) w ramach schematu ekspansji półgrupowej. Badania te były możliwe dzięki połączeniu teoretycznych wniosków i wspomaganych komputerowo obliczeń w Pythonie. Kluczowym aspektem tej pracy było podkreślenie "warunku rezonansowego" — matematycznego kryterium zapewniającego zgodność z wyjściową strukturą algebry AdS. Dodatkowo wprowadziliśmy element pochłaniający $0_S$, przypisując jednocześnie elementowi $s_0$ (powiązanemu z generatorem Lorentza $J_{ab}$) rolę elementu tożsamości. Matematycznie oznacza to przejście od półgrup do monoidów, podczas gdy warunek rezonansowy obejmuje jedynie warunek parzystości ($s_{even}\cdot s_{even}=s_{even},\, s_{even}\cdot s_{odd} = s_{odd},\, s_{odd}\cdot s_{odd}=s_{even}$, tj. elementy oznaczone jako parzyste przez "mnożenie" dają parzysty element itd.).

Zaczynając od brute-forcowego wyszukiwania monoidów (półgrup z elementem neutralnym) z warunkiem parzystości, mogliśmy wyprowadzić relacje komutacyjne dla nowych generatorów. Kluczowym osiągnięciem była kompleksowa klasyfikacja rezonansowych algebr oparta na monoidach z parzystością, wykraczająca poza wcześniejsze prace (\textbf{H3}). Te rozszerzone ramy ujawniły zróżnicowany wachlarz możliwości algebraicznych, z których każdy teoretycznie odpowiada różnym Lagrangianom i równaniom pola, z potencjalnymi zastosowaniami w przyszłych badaniach.

%%%%%%%%%%%%%%%%%%%%%%%%%%%%%%%%%%%%%%%%%%%%%%%%%%%%%%%%%%%%%%%%%%%%%%%%%%%%
\myparagraph{Resonant Superalgebras and $\mathcal{N}=1$ Supergravity (H6)}

Opierając się na wcześniejszych pracach dotyczących bozonowych rezonansowych algebr (\textbf{H2, H3, H4, H5}), niniejsze studium przenosi się w obszar rozszerzeń supersymetrycznych $\mathcal{N}=1$ z jednym generatorem superładunku. Głównym celem (\textbf{H6}) jest odpowiedź na pytanie, ile prawidłowych wersji supersymetrycznych rezonansowych algebr można znaleźć dla zawartości generatorów $JPZ+Q$. Aby znaleźć wszystkie możliwe realizacje, musieliśmy najpierw zapytać, jakie cechy są niezbędne do uzyskania "dobrej" superalgebry oraz co jest wymagane, aby skutecznie zbudować teorię supergrawitacji.

Rozległe poszukiwania wszystkich możliwości, wspomagane systemem symbolicznej algebry komputerowej \textit{Cadabra}, ostatecznie doprowadziły do znalezienia wszystkich superalgebr $JPZ+Q$, z których sześć prowadzi do sformułowania supergrawitacji, podczas gdy pozostałe trzy reprezentują przypadki niestandardowe. Wśród nich dwie stanowią już znane rozszerzenia supersymetryczne algebr $\mathfrak{C}_{4}$ i $\mathfrak{B}_{4}$, odpowiednio Soroka-Soroka \cite{Soroka:2006aj} oraz Maxwell \cite{Lukierski:2010dy}:
\begin{equation*}
	\begin{array}{llllllll}
		\mathfrak{C}_{4}: &  &
		\begin{tabular}[t]{c|ccc}
			\lbrack .,.] & J & P & Z \\ \hline
			J & J & P & Z \\
			P & P & Z & P \\
			Z & Z & P & Z
		\end{tabular}
		&  &
		\begin{tabular}[t]{c|c}
			\lbrack .,.] & Q \\ \hline
			J & Q \\
			P & Q \\
			Z & Q
		\end{tabular}
		&  &
		\begin{tabular}[t]{c|c}
			\{.,.\} & Q \\ \hline
			Q & P+Z
		\end{tabular}
		& ~~~\\[55pt]
		\mathfrak{B}_{4}: &  &
		\begin{tabular}[t]{c|ccc}
			\lbrack .,.] & J & P & Z \\ \hline
			J & J & P & Z \\
			P & P & Z & 0 \\
			Z & Z & 0 & 0
		\end{tabular}
		&  &
		\begin{tabular}[t]{c|c}
			\lbrack .,.] & Q \\ \hline
			J & Q \\
			P & 0 \\
			Z & 0
		\end{tabular}
		&  &
		\begin{tabular}[t]{c|c}
			\{.,.\} & Q \\ \hline
			Q & Z
		\end{tabular}
		& ~~~
	\end{array}
\end{equation*}
Z drugiej strony, wszystkie nowe supersymetryczne rozszerzenia algebr rezonansowych mają podstrukturę podobną do Poincar\'e- \textbf{(H3, H4, H5)}. Dla porównania, poniżej przedstawiamy tylko trzy z nowych przykładów:
\begin{equation*}
	\begin{array}{llllllll}
		C_{4}: &  &
		\begin{tabular}[t]{c|ccc}
			\lbrack .,.] & J & P & Z \\ \hline
			J & J & P & Z \\
			P & P & 0 & P \\
			Z & Z & P & Z
		\end{tabular}
		&  &
		\begin{tabular}[t]{c|c}
			\lbrack .,.] & Q \\ \hline
			J & Q \\
			P & 0 \\
			Z & Q
		\end{tabular}
		&  &
		\begin{tabular}[t]{c|c}
			\{.,.\} & Q \\ \hline
			Q & P
		\end{tabular}
		&  \\[55pt]
		\tilde{C}_{4}: &  &
		\begin{tabular}[t]{c|ccc}
			\lbrack .,.] & J & P & Z \\ \hline
			J & J & P & Z \\
			P & P & 0 & 0 \\
			Z & Z & 0 & Z%
		\end{tabular}
		&  &
		\begin{tabular}[t]{c|c}
			\lbrack .,.] & Q \\ \hline
			J & Q \\
			P & 0 \\
			Z & 0
		\end{tabular}
		&  &
		\begin{tabular}[t]{c|c}
			\{.,.\} & Q \\ \hline
			Q & P
		\end{tabular}
		& $\textcircled{\footnotesize 1}$ \\[55pt]
		\tilde{C}_{4}: &  &
		\begin{tabular}[t]{c|ccc}
			\lbrack .,.] & J & P & Z \\ \hline
			J & J & P & Z \\
			P & P & 0 & 0 \\
			Z & Z & 0 & Z%
		\end{tabular}
		&  &
		\begin{tabular}[t]{c|c}
			\lbrack .,.] & Q \\ \hline
			J & Q \\
			P & 0 \\
			Z & Q
		\end{tabular}
		&  &
		\begin{tabular}[t]{c|c}
			\{.,.\} & Q \\ \hline
			Q & Z
		\end{tabular}
		& $\textcircled{\footnotesize 2}$ 
	\end{array}
\end{equation*}
Należy zauważyć, że dobrze ugruntowane schematy supersymetryczne (Poincaré i AdS) dla $JP+Q$ zostały już przedstawione we wstępie pod \eqref{table_Poincare} i \eqref{table_AdS}.

Nasze wysiłki tutaj zaowocowały sformułowaniem działań Cherna-Simonsa, które są niezmiennicze względem odpowiednich rezonansowych superalgebr. Odkrycia te odgrywają kluczową rolę we włączaniu superalgebr do szerszych ram teorii supergrawitacyjnych, tym samym pogłębiając nasze zrozumienie złożonych relacji między różnymi teoriami supergrawitacji i supersymetrycznymi działaniami CS opartymi na odmiennych rezonansowych supersymetriach.

Niniejsza praca systematycznie kategoryzuje wszystkie możliwe  $\mathcal{N}=1$ rozszerzenia rezonansowych speralgebr Spośród nich sześć rozszerzeń prowadzi do spójnych działań CS supergrawitacji, podczas gdy pozostałe rozszerzenia są klasyfikowane jako niestandardowe, charakteryzujące się brakiem członu translacyjnego w antykomutatorze $\{Q, Q\}$. 

Charakterystyka i implikacje tych działań w $3D$ zostały szczegółowo omówione wcześniej w kontekście własności \eqref{Models}. Podsumowanie ilościowe algebr zostało również zaprezentowane we wstępie \eqref{Overview}.

Kluczowym aspektem badania było wykorzystanie symbolicznego systemu algebry komputerowej \textit{Cadabra}, który wspomógł rozległe poszukiwania potencjalnych superalgebr. Obliczenia, wymagające 27 godzin czasu pracy na standardowym komputerze PC, umożliwiły kompleksową analizę super-tożsamości Jacobiego. Wykazały one również znaczącą barierę obliczeniową dla \textit{Cadabry}, ponieważ podobna analiza dla $\mathcal{N}=2$ wymagałaby od jednego do dwóch rzędów wielkości więcej zasobów obliczeniowych. Ten problem został rozwiązany i pomyślnie przezwyciężony w kolejnej publikacji (\textbf{H9}). 

Warto zauważyć, że wyprowadzenie supersymetrycznych rozszerzeń rezonansowych algebr w tej pracy opierało się na czysto brute-force analizie super-tożsamości Jacobiego, co stanowiło odejście od powszechnie stosowanej metody rozszerzeń półgrupowych. To metodologiczne odstępstwo było kluczowe w odkryciu nowych rozszerzeń i przygotowało grunt pod moje dalsze badania opisane w (\textbf{H9, H10, H12}).

%%%%%%%%%%%%%%%%%%%%%%%%%%%%%%%%%%%%%%%%%%%%%%%%%%%%%%%%%%%%%%%%%%%%%%%%%%%%
\myparagraph{Resonant Superalgebras for Supergravity (H9)}

Omawiając historyczny kontekst algebr Maxwella \cite{Bacry:1970ye, Schrader:1972zd} i Soroki-Soroki \cite{Soroka:2006aj}, a także S-rozszerzonych algebr, praca podkreśla szersze uogólnienie osiągnięte poprzez rezonansowe (super)algebry. Struktury te dziedziczą bazowy wzorzec (anty)-komutacji AdS i można je traktować jako algebraiczne rozszerzenia przypadków Poincaré lub AdS.

We współpracy z dr. K. Graczykiem w pracy (\textbf{H9}) udało mi się opracować kompleksowy katalog rezonansowych (super)algebr (już wspomniany w sekcji \eqref{Overview}), obejmujący setki przykładów odpowiednich do konstruowania modeli supergrawitacji, z których każdy posiada unikalną realizację Lagrangianu. Było to możliwe dzięki rozwiązaniu złożoności obliczeniowych oraz zastosowaniu technik stochastycznych. Opracowana wydajna metodologia, oparta na rozpoznawaniu wzorców, znacząco usprawniła wcześniejsze podejścia, rozszerzając standardowy schemat $JP+Q$ do bardziej złożonych konfiguracji, obejmujących nawet podwojenie generatorów typu Lorentza, translacyjnych i supersymetrycznych.

Szablon rezonansowy z poprzedniej sekcji \eqref{empty}, bez wypełnienia pustych miejsc, służy jedynie jako wskazówka do poszukiwań. Postęp w identyfikacji przykładów spełniających wszystkie warunki opierał się na kilku ręcznie opracowanych przypadkach bazujących na intuicji. To skłoniło mnie do rozważenia różnych metod komputerowych. Samodzielnie mogłem przeprowadzać obliczenia tylko dla ograniczonej liczby generatorów (\textbf{H3, H5, H6}). Podejście brute-force zastosowane do bardziej ogólnego układu wiązałoby się z nieosiągalnie dużą przestrzenią poszukiwań. W przypadku $JPZU+QY$ nieuniknione byłoby sprawdzenie $35$ unikalnych supertożsamości Jacobiego z sześcioma (super)komutatorami do wykonania dla każdego z $344\,373\,768$ możliwych kandydatów na superalgebrę.

Aby poradzić sobie z tym problemem, zaproponowaliśmy nową metodologię, która opiera się na wykorzystaniu zaawansowanych metod komputerowych do generowania kandydatów spełniających warunki rezonansowe, kierując się zadanymi ograniczeniami w celu wyeliminowania niefizycznych przypadków.

Nasza metoda może zostać podsumowana w trzech etapach:
\begin{itemize}
	\item (a) „inteligentne” generowanie kandydatów algebraicznych,
	\item (b) użycie wyszukiwania wzorców zamiast pełnej ewaluacji tożsamości Jacobiego,
	\item (c) zastosowanie algorytmu stochastycznego, który „uczy się” rodzaju relacji Jacobiego najbardziej krytycznych do sprawdzenia.
\end{itemize}

W pierwszym kroku generujemy możliwych kandydatów, z których każdy reprezentuje unikalną konfigurację generatorów spełniającą warunki rezonansowe i zapisaną w postaci tablic algebry. Od samego początku wprowadzamy kilka ograniczeń odrzucających oczywiste niepoprawne lub niefizyczne sytuacje.

Algebra jest poprawnie zdefiniowana, gdy wszystkie tożsamości Jacobiego są spełnione\\[-15pt]
\begin{align}
	\label{Eq:Jacobi}
	0 \equiv &\left( \left(X_i, X_j\right), X_k\right)\oplus \left(\left( X_j, X_k\right), X_i \right) \oplus \left(\left( X_k, X_i\right), X_j\right)\,,
\end{align}
gdzie $\left( .\,,.\right)$ odnosi się do $\{.\,,.\}$ lub $[. \, , .]$; znak $\oplus$ uwzględnia możliwy czynnik $(-1)$ w gradacyjnej relacji Jacobiego, a $X_{i}, X_{j}, X_{k}$ odnosi się do dowolnych generatorów (super)algebry.
%
Główna idea naszego podejścia opartego na wzorcach (zamiast bezpośredniej ewaluacji)
\begin{equation}
	\label{Eq:Pattern}
	Pattern = { \left( \left(X_i, X_j\right), X_k\right),\,\left(\left( X_j, X_k\right), X_i \right) ,\, \left(\left( X_k, X_i\right), X_j\right) }
\end{equation}
opiera się na obserwacji, że \eqref{Eq:Jacobi} zeruje się z powodu zakodowanych relacji między stałymi struktury AdS. W efekcie spodziewamy się trzech instancji tego samego obiektu, w $Pattern = (L, L, L)$, gdzie $L$ odnosi się do $0$, $X_i$ lub kombinacji liniowej $X_i$.

Aby dodatkowo przyspieszyć proces, zastosowaliśmy algorytm, który \textit{„uczy się”}, które typy tożsamości Jacobiego (z długiej listy wszystkich trójek generatorów) są najbardziej krytyczne, tzn. te, dla których test Jacobiego zawodzi. Dalsza weryfikacja rozpoczyna się od sprawdzania krytycznego przypadku. Warto zauważyć, że proces ten jest dynamiczny, a krytyczne przypadki są identyfikowane w trakcie całego procesu.

Jak pokazano na Rys.~\ref{fig:disttributions} z histogramem krytycznych relacji, pojedyncze sprawdzenie $\mathrm{Jacobi}(P, Y, Q)$ natychmiast eliminowało ok. 100\,000 błędnych kandydatów. Dla bardziej złożonych przypadków jak $\mathcal{N}=2$ supergrawitacja z dodatkowym $U(1)$ generatorem $T$, metodologia ta może być dostosowana do bardziej złożonych wzorców (\textbf{H10}).

\begin{figure*}[h!]
	\centering
	\includegraphics[width=0.78\textwidth]{histogram.PNG}
	\caption{Ostateczne rozkłady najbardziej „problematycznych” tożsamości Jacobiego podczas procesu wyszukiwania. „PYQ” odnosi się do $\mathrm{Jacobi}(P,Y,Q)$ itd. Źródło: (\textbf{H9})}
	\label{fig:disttributions}
\end{figure*}

Ostatni analizowany scenariusz dotyczył przykładów domkniętych superalgebr z efektywnie dwiema grupami generatorów podobnych do Lorentza, translacyjnych i superładunkowych. Te supersymetryczne rozszerzenia, wywodzące się z trzydziestu bozonowych rodzin algebraicznych $JPZU$, wymagały apriori ponad $10^{10}$ sprawdzeń (osiągniętych w ciągu \textit{24 godzin} zamiast szacowanych \textit{dwóch miesięcy} dedykowanych obliczeń symbolicznych w oprogramowaniu \textit{Cadabra}). 

Niniejsza praca wnosi istotny wkład w teorię supergrawitacji, rozszerzając struktury algebraiczne poza dotychczasowe granice. Dzięki wprowadzeniu nowych generatorów bozonowych i fermionowych oraz zastosowaniu wydajnej metody obliczeniowej, otwiera nowe możliwości budowy unikalnych modeli supergrawitacji.

Podejście oparte na wzorcach w klasyfikacji rezonansowych (super)algebr pokazuje, jak zaawansowane narzędzia obliczeniowe mogą radzić sobie z rozległymi i złożonymi przestrzeniami wyszukiwania, prowadząc do pełnego katalogu możliwości. Wiele elementów sekcji dotyczących konstrukcji \eqref{Construction} oraz przeglądu \eqref{Overview} rezonansowych (super)algebr pochodzi bezpośrednio z tej pracy (\textbf{H9}), co ukazuje jej kluczową rolę w rozwoju rezonansowych superalgebr i ich zastosowań w supergrawitacji.

%%%%%%%%%%%%%%%%%%%%%%%%%%%%%%%%%%%%%%%%%%%%%%%%%%%%%%%%%%%%%%%%%%
\myparagraph{N=2 Resonant Superalgebra for Supergravity (H10)} 

Teoria opracowana w poprzedniej publikacji (\textbf{H9}), mimo wprowadzenia dwóch superładunków, nie uwzględniała rozszerzenia o centralny ładunek \cite{Haag:1974qh}, oferującego generator $U(1)$ jako wynik antykomutacji $Q$ z $Y$. Dlatego dwa fermionowe ładunki pojawiają się podobnie jak w pracach \cite{Durka:2012wd, Soroka:2006aj, Lukierski:2010dy, Durka:2011gm}. 

Po pracy (\textbf{H6}) z $\mathcal{N}=1$ teoriami supergrawitacji dla $JPZ+Q$, (\textbf{H10}) bada $\mathcal{N} = 2$, $D = 3, 4$ scenariusze wprowadzające $A_{\mu}$, pole cechowania $U(1)$ \cite{Gocanin:2019ota}. Standardowa supergrawitacja $\mathcal{N}=2$, w szczególności, łączy dynamikę grawitonu i grawitina z dodatkowymi grawitinem oraz wektorowym polem spinowym-1, wynikającymi z rozszerzonej grupy supersymetrii, która uwzględnia centralne ładunki i wewnętrzne rotacje dwóch grawitin \cite{VanNieuwenhuizen:1981ae}.

Należy zauważyć, że typowy komutator $[P_a, Q^I] \sim Q^I$ lub $0$, jednak w (\textbf{H9}) opisaliśmy dziesięć konfiguracji $JP+QY$ (z dwoma superładunkami, ale bez generatora $T$), gdzie $[P_a, Q^I] \sim Q^{J}$. Takie algebry, wg naszej wiedzy, nie były badane w literaturze, gdyż dla dwóch superładunków natychmiast przechodzi się do konfiguracji z centralnym ładunkiem.

Analiza konfiguracji $JPZ+QY+T$ ostatecznie doprowadziła nas do odkrycia pojedynczej realizacji algebraicznej, zakorzenionej w algebrze Soroka-Soroka, która obejmuje generatory Lorentza $J_{ab}$, generatory Lorentzowsko-podobne $Z_{ab}$, translacje $P_a$ oraz generator $T$ wraz z dwoma superładunkami $Q^I$, będącymi $Q$ i $Y$. Ta superalgebra reprezentuje jedyną możliwą realizację dla tego zestawu generatorów.

Odpowiadająca konstrukcja $3D$ supergrawitacji Cherna-Simonsa oparta na tej nowej superalgebrze ponownie zapewnia pełną niezmienniczość symetrii, dualne definicje generatorów oraz niezbędne elementy do skonstruowania Lagranżjanu: koneksji, superkrzywizny, parametrów cechowania i form Killinga. W artykule przedyskutowano możliwość rozszerzenia tego podejścia na konstrukcję czterowymiarową ($4D$), podkreślając szerszą perspektywę domykania rezonansowych algebr.

%%%%%%%%%%%%%%%%%%%%%%%%%%%%%%%%%%%%%%%%%%%%%%%%%%%%%%%%%%%%%%%%%%%%%%%%%%%%
\myparagraph{Resonant enlargements from pattern-based analysis (H12)} 

Praca ta stanowi ponownie przegląd ostatnich osiągnięć zaprezentowanych w (\textbf{H9}) i (\textbf{H10}), dostarczając zwięzłego podsumowania zastosowanych metod oraz wykorzystania techniki wyszukiwania opartego na wzorcach w konstruowaniu rezonansowych (super)algebr. Konsoliduje kluczowe spostrzeżenia dotyczące formułowania nowych struktur algebraicznych, podkreślając ich znaczenie dla teorii supergrawitacji. 

Motywacja do badania rezonansowych (super)algebr jako uogólnień (super)algebr Poincarégo i AdS została przedstawiona w analogii do fermionowych teorii $\mathcal{N}$ rozszerzonych, koncentrując się teraz także na rozszerzeniu generatorów bozonowych.

Praca wprowadza działanie Cherna-Simonsa dla grawitacji w trzech wymiarach ($3D$) i podkreśla, w jaki sposób supersymetria rozszerza ten schemat do supergrawitacji. Motywacja do badania rezonansowych (super)algebr, jako uogólnień (super)algebr Poincarégo i AdS, została przedstawiona w analogii do teorii rozszerzonych $\mathcal{N}$, koncentrując się na rozszerzaniu generatorów bozonowych.

Zastosowanie efektywnej metody obliczeniowej opartej na wzorcach do generowania rezonansowych struktur doprowadziło do odkrycia szerokiej klasy nowych (super)algebr Liego, rozszerzając zakres możliwości dla teorii fizycznych. Algorytm upraszcza ewaluację tożsamości super-Jacobiego poprzez dynamiczne identyfikowanie kluczowych przypadków, co znacząco redukuje złożoność procesu wyszukiwania. Parafrazując:  algorytm „uczy się” najbardziej problematycznych tożsamości Jacobiego i wykorzystuje je w procesie wyszukiwania. Zwróciliśmy również uwagę na wyzwania związane z rozszerzeniem całej struktury na $\mathcal{N}=2$ oraz dodatkowe wymagania do uwzględnienia dodatkowych generatorów.

Na koniec omówiono znaczenie algebry rezonansowej dla badania nowych modeli grawitacji oraz potencjalne kierunki przyszłych badań. Praca ta podkreśla istotność tych osiągnięć w poszerzaniu naszego zrozumienia (super)grawitacji i jej rozszerzeń.

W artykule wspominamy również o jednej interesującej możliwości, mianowicie \textit{bi-supergrawitacji}. Drugi fermionowy superładunek oferuje scenariusz bi-supergrawitacji, analogiczny do teorii bi-metrycznej. Dla czysto bozonowej teorii, zamiast arbitralnego członu interakcji w działaniu bi-metrycznym (zob. Ref.~\cite{Durka:2011va, Hoseinzadeh:2017lkh}), możemy spróbować dołączyć do metryki bazowej (przypisanej do $J_{ab}, P_{a}$) inne pole tensorowe rzędu 2 (odpowiadającemu parze $\hat{J}_{ab},\hat{P}_{a}$). Mając dwie kopie generatorów Lorentzowsko-podobnych i translacyjnie-podobnych, otrzymujemy $30$ przypadków bozonowych algebr rezonansowych (\textbf{H3, H5}). Jeśli $J_{ab},P_a \to (\omega^{ab}, e^a ) \to g_{\mu\nu}$, możemy rozważyć $\hat{J}_{ab}, \hat{P}_a \to (h^{ab}, k^a) \to f_{\mu\nu}$ jako inne pole tensorowe rzędu drugiego sprzężone z grawitacją. Zamiast arbitralnie mieszanych członów, mamy $30$ teorii \textit{bi-koneksyjnych} z odrębnymi Lagrangianami wynikającymi z odpowiadających $30$ algebr:
\begin{itemize}
	\item $e^a$ i $\omega^{ab}$ (vierbein i koneksja spinowa), opisujące grawitację z odpowiadającymi generatorami $P_a$ i $J_{ab}$,
	\item $\hat{e}^a$ i $\hat{\omega}^{ab}$, nowe pola z odpowiadającymi dodatkowymi generatorami $\hat{P}_a$ i $\hat{J}_{ab}$.
\end{itemize}

Jeśli pójść dalej i uogólnić układ do $JPZU+QY$ (efektywnie podwajając $2\times (JP+Q)$), można by było otrzymać bi-supergrawitację, gdzie standardowe teorie $\mathcal{N}=2$ mogły by być postrzegane jako pewne złamanie symetrii lub granica. Konstrukcyjnie w $3D$ można podążać za konkretną realizacją superalgebry w ramach konstrukcji Cherna-Simonsa \eqref{CS}, natomiast dla $4D$ można zastosować działanie typu Borna-Infelda \cite{Concha:2014vka}, model MacDowella-Mansouriego \cite{MacDowell:1977jt} lub zdeformowaną teorię BF \cite{Durka:2011nf, Durka:2012wd}.

%%%%%%%%%%%%%%%%%%%%%%%%%%%%%%%%%%%%%%%%%%%%%%%%%%%%%%%%%%%%%%%%%%%%%%%%%%%%%%%%%

\subsubsection{Modele grawitacyjne i symetrie (H1, H7, H8, H11)}

Moje badania koncentrują się również na roli innych symetrii w formułowaniu i analizie modeli grawitacyjnych, w tym na eksploracji różnych modeli, ich implikacjach dla zachowań asymptotycznych oraz ich rozwiązaniach.
\begin{itemize}
	\item \textit{Struktury brzegowe w grawitacji} (\textbf{H7}): Zbadałem struktury brzegowe w teorii grawitacji sformułowanej jako teoria BF, szczególnie w kontekście członów topologicznych i ich wpływu na ładunki asymptotyczne. Praca ta wyjaśnia różnicę między symetriami w objętości a symetriami na brzegu oraz fizyczne konsekwencje tych różnic.
	
	\item \textit{Porównanie dwóch teorii BF} (\textbf{H11}): Porównałem dwa różne sformułowania BF grawitacji (typu MacDowell-Mansouri i typu Plebańskiego), analizując ich asymptotyczne ładunki i metody regularyzacji. Badanie to wskazuje na lepsze radzenie sobie z zachowaniem asymptotycznym przez teorię MacDowella-Mansouriego i jej odporność w różnych kontekstach grawitacyjnych (w przestrzeniach asymptotycznie AdS nie są wymagane dodatkowe regularyzacyjne człony, w przeciwieństwie do teorii Plebańskiego). Praca oferuje szczegółowe wyprowadzenie zachowanych ładunków "corner'owych" (narożnych) związanych z lokalnymi symetriami cechowania w obu teoriach i wykorzystuje te wyrażenia do obliczania ładunków asymptotycznych dla różnych rozwiązań, takich jak AdS-Schwarzschild czy AdS-Taub-NUT.
	
	\item \textit{Rotacja w czasoprzestrzeni Taub-NUT} (\textbf{H8}): Zaproponowałem nową interpretację parametru NUT w czasoprzestrzeni Taub-NUT, tradycyjnie postrzeganego jako masa dualna będąca grawitacyjnym odpowiednikiem monopolu magnetycznego. Nowa interpretacja rozwiązuje szereg niespójności w istniejącej literaturze, wykazując, że parametr ten reprezentuje nietrywialną formę rotacji związaną z strunami Misnera. Parametr $n$ przyczynia się zarówno do momentu pędu (będącego ładunkiem Noether w nieskończoności, obliczonym dla symetrii dyfeomorfizmu odpowiadającej obrotowemu wektorowi Killinga, z kluczową rolą członów Holsta i Pontryagina dodanych do teorii Einsteina), jak i do entropii horyzontu w termodynamice czarnych dziur.
	
	\item \textit{Stopnie swobody w grawitacji Pure Lovelocka} (\textbf{H1}): Przeanalizowałem kanoniczną strukturę tzw. grawitacji Pure Lovelock, gdzie działanie składa się ze stałej kosmologicznej i pojedynczego wielomianu wyższego rzędu w tensorze Riemanna. Jest to wyższejwymiarowe uogólnienie Ogólnej Teorii Względności, ujawniające unikalne właściwości dotyczące liczby fizycznych stopni swobody i ich zależności od geometrii tła. Praca ta dostarcza wglądu w różnice i podobieństwa między grawitacją Lovelocka a standardowym modelem OTW.
	
\end{itemize}

%%%%%%%%%%%%%%%%%%%%%%%%%%%%%%%%%%%%%%%%%%%%%%%%%%%%%%%%%%%%%%%%%%%%%%%%%%%%
\myparagraph{Corners of gravity: the case of gravity as BF theory (H7)} 

Praca oferuje szczegółowe badanie struktury brzegowej w kontekście grawitacji sformułowanej jako teoria BF typu MacDowella-Mansouriego \cite{Durka:2012wd, Freidel:2005ak} ze względu na jej ogólność, która pozwala na uwzględnienie stałej kosmologicznej oraz eksplorację wpływu członów topologicznych (Eulera, Pontryagina, Nieh-Yana) obecnych w działaniu tej teorii.  W czasoprzestrzeni bez brzegów niezmienniki topologiczne nie odgrywają żadnej roli, ale w przypadku obecności brzegów wpływają na brzegową strukturę symplektyczną i ładunki, czyniąc teorię inną od tych omawianych w literaturze \cite{Freidel:2020xyx}.

Praca wpisuje się w badania nad symetriami grawitacyjnymi i ich związkiem z zachowanymi ładunkami. W artykule pokazano, że ustalone formuły na ładunki są zależne od stałej kosmologicznej oraz członów topologicznych, zapewniając lepsze zrozumienie tych aspektów w teoriach grawitacji. Na przykładzie rozwiązania AdS-Schwarzschild pokazano, że ładunki bezpośrednio obliczone z działania BF dają poprawne wartości. Uwzględnienie członów topologicznych regularyzuje wkłady brzegowe, unikając rozbieżności w zachowanych wielkościach, takich jak masa (oraz moment pędu, jeśli rozważa się czasoprzestrzeń Kerra).

Badanie to wnosi wkład w trwające wysiłki związane z regularyzacją i interpretacją rozwiązań. Podkreśla ono różnicę między symetriami cechowania (AdS obejmującą zarówno: Lorentza, jak i translacje) w objętości czasoprzestrzennej a symetriami fizycznymi na brzegu. Ponadto, zagłębiamy się w teorię Einsteina-Cartana-Holsta oraz związaną z nią strukturę symplektyczną, dostarczając szczegółowego badania potencjałów symplektycznych na powierzchni $\Sigma$ oraz w cornerze $S$. Termin „corner” odnosi się do specyficznego pojęcia w ramach struktur brzegowych układu grawitacyjnego. Tradycyjne podejścia do grawitacji często koncentrują się na właściwościach objętości czasoprzestrzennej, jednak niedawne wyniki uwypukliły znaczenie brzegów i ich ładunków. „Corner” reprezentuje bardziej wyrafinowaną strukturę brzegową.

Obliczamy ładunki, podstawiając transformacje Lorentza, translacje i przekształcenia dyfeomorficzne do zdekomponowanej struktury symplektycznej wyprowadzonej z potencjałów symplektycznych. Analizując bliżej ładunki związane z lokalną symetrią Lorentza oraz dyfeomorfizmami zauważyć możemy, że translacyjny składnik ładunku $SO(3,2)$ znika zarówno na powierzchni $\Sigma$, jak i w cornerze $S$, wskazując, że jest on czystym cechowaniem. Prowadzi to do uproszczenia formuł i lepszego zrozumienia fizycznych implikacji tych symetrii. Symetrie cechowania w objętości, takie jak lokalne transformacje Lorentza i translacje, są zestawiane z fizycznymi symetriami pojawiającymi się na brzegach, podkreślając różnicę między swobodą cechowania w objętości a fizycznymi ładunkami na brzegu.

%%%%%%%%%%%%%%%%%%%%%%%%%%%%%%%%%%%%%%%%%%%%%%%%%%%%%%%%%%%%%%%%%%%%%%%%%%%
\myparagraph{A tale of two theories of gravity in asymptotically AdS (H11)} 

Kolejny artykuł definiuje dwie różne grawitacyjne teorie BF, omawiając ich podobieństwa i różnice, a także ukazując ich unikalne cechy oraz konsekwencje dla badań nad grawitacją. Możnemy tutaj podkreślić efektywność teorii BF typu MacDowella-Mansouriego \cite{Freidel:2005ak, Durka:2012wd} w porównaniu do teorii BF typu Plebańskiego \cite{Plebanski:1977zz} w kontekście opisu ładunków asymptotycznych i unikania potrzeby stosowania regularyzacji kontrczłonami, co zapewnia bardziej spójne i kompletne ramy teoretyczne.

Rzeczywiście, pierwsze sformułowanie ma kilka zalet w porównaniu ze standardowym sformułowaniem teorii Einsteina-Hilberta-Holsta, w którym to grawitacja jest zdefiniowana jako teoria BF $SO(3,1)$ lub $SO(2,3)$ z dodanym parametrem Barbero-Immirziego $\gamma$ i nałożonym więzem $B = (\gamma^{-1} +\star)\, e\wedge e$, które jest teraz skutecznie zastępione przez warunek gdzie pole $B$ jest proporcjonalne do krzywizny Anty de Sittera. Po pierwsze, teoria MacDowella-Mansouriego jest jawnie lokalnie niezmiennicza względem transformacji Lorentza na każdym etapie obliczeń struktury symplektycznej brzegu i ładunków, a także nie wymaga problematycznego narzucania "simplicity constraint". Po drugie, zawiera stałą kosmologiczną jako niezbędny składnik od samego początku, co czyni ją naturalnym narzędziem do badania czasoprzestrzeni asymptotycznie AdS oraz związanych z nimi ładunków.

Artykuł wpisuje się w szerszy nurt badań nad grawitacją, który koncentruje się coraz bardziej na brzegach i skończonych asymptotycznych regionach czasoprzestrzeni, szczególnie w kontekście symetrii asymptotycznych. Ten kierunek badawczy doprowadził do wzrostu zainteresowania ładunkami Noether związanymi z tymi brzegami. Ładunki w teorii grawitacji opisanej działaniem Einsteina–Cartana–Holsta z dodatkowymi członami topologicznymi były omawiane w przeszłości w \cite{Durka:2012wd, Godazgar:2020gqd, Godazgar:2020kqd}. Badania (\textbf{H11}) dotyczą zachowanych ładunków powierzchniowych związanych z lokalnymi symetriami cechowania w obu teoriach BF, które są niezmiennicze względem lokalnych transformacji Lorentza i dyfeomorfizmach. Ładunki są zdefiniowane za pomocą wzoru, który wiąże je z formą symplektyczną, zapewniając fundamentalną podstawę do ich obliczenia. 

Wykonując jawne obliczenia ładunków asymptotycznych w obu teoriach dla różnych rozwiązań, takich jak AdS-Schwarzschild, AdS-Kerr i AdS-Taub-NUT, okazuje się, że teoria Plebańskiego prowadzi do rozbieżnych ładunków, podczas gdy w teorii MacDowella-Mansouriego ładunki są skończone. W przypadku rozwiązania AdS-Taub-NUT badanie pokazało, że ładunki asymptotyczne zależą od parametru Barbero-Immirziego, w przeciwieństwie do wcześniej omawianych rozwiązań.

%%%%%%%%%%%%%%%%%%%%%%%%%%%%%%%%%%%%%%%%%%%%%%%%%%%%%%%%%%%%%%%%%%%%%%%%%%%%

\myparagraph{The first law of black hole thermodynamics for Taub–NUT (H8)} 

Niniejszy artykuł redefiniuje rozumienie czasoprzestrzeni Taub-NUT, przedstawiając nowe spojrzenie na naturę parametru NUT $n$ i jego implikacje dla termodynamiki czarnych dziur i ładunków Noether. Tradycyjnie uważa się, że czasoprzestrzeń Taub-NUT posiada dualny ładunek do masy, będący grawitacyjnym analogiem monopolu magnetycznego. Niniejsze badanie pokazuje natomiast, że $n$ prowadzi do unikalnej formy rotacji wynikającej ze strun Misnera. Struny te, wcześniej odrzucane jako artefakty współrzędnych, okazują się kluczowymi wkładami do momentu pędu, $J\sim n M$ będącego ładunkiem Noether w nieskończoności (\textbf{H11}) (obliczonym dla symetrii dyfeomorficznej) i entropii (odpowiednio wykorzystując definicję horyzontu Killinga). Moje badanie wprowadziło nową interpretację, w której przeciwnie obracające się dwie struny Misnera fundamentalnie redefiniują dynamikę momentu pędu całego układu. Całkowanie w trakcie obliczania ładunków doprowadziło wielu badaczy do przeoczenia tego kluczowego zjawiska. Teraz można zauważyć, że zanik całkowitego momentu pędu wynika z wzajemnego znoszenia się niezerowych momentów pędu o przeciwnych znakach w dwóch przeciwległych półkulach.

Moja praca dotyczy zarówno obliczeniowych, jak i koncepcyjnych wyzwań w zrozumieniu czasoprzestrzeni Taub-NUT, podkreślając ograniczenia standardowej metody euklidezacji. Metody te nadmiernie upraszczają system - prowadząc wciaż do niespójności poprzez pominięcie istotnych cech reżimu lorentzowskiego. W przeciwieństwie do tego, niniejsze badanie obejmuje pełną złożoność czasoprzestrzeni lorentzowskiej, wraz z koniecznością uwzględnienia członów topologicznych w ramach Einsteina-Cartana (ponownie wykazując dużą użyteczność teorii BF) w celu wyprowadzenia skorygowanych relacji termodynamicznych.

Ta nowa interpretacja nie tylko rozwiązuje kilka niejasnych aspektów występujących w istniejącej literaturze, aale także ustanawia spójne ramy dla opisu termodynamiki czarnych dziur w lorentzowskiej czasoprzestrzeni Taub-NUT. W badaniu ponownie przeanalizowano pierwsze prawo termodynamiki czarnych dziur, dostosowując je tak, aby wyraźnie uwzględniało wkład strun Misnera. Prawo to zawiera teraz wyrażenie, które odzwierciedla efekty rotacyjne i entropię związaną ze strunami, wyrażone jako:
\begin{align}
	dM = T dS + \Omega dJ - \frac{1}{2n} d\left(\frac{nr_+}{G}\right).
\end{align}
To sformułowanie, które otrzymałem jako jeden z pierwszych, ujawnia skomplikowaną równowagę między termodynamicznymi właściwościami horyzontu a wkładem strun Misnera, rzucając światło na wcześniej niezbadane aspekty struktury czasoprzestrzeni.

%%%%%%%%%%%%%%%%%%%%%%%%%%%%%%%%%%%%%%%%%%%%%%%%%%%%%%%%%%%%%%%%%%%%%%%%%%%%
\myparagraph{Dynamical Structure of Pure Lovelock Gravity (H1)}

Na podstawie mojego wcześniejszego doświadczenia \cite{Durka:2012wd, Durka:2010zx}, przeprowadziłem tutaj analizę kanoniczną dynamicznych cech grawitacji Pure Lovelock (PL) dla czasoprzestrzeni o wymiarach większych niż cztery. Teoria ta zawiera wielomian $p$-tego rzędu w tensorze Riemanna w $(d+1)$-wymiarach jako człon kinetyczny (bez liniowego członu w skalarnej krzywiźnie), zachowując równania ruchu drugiego rzędu względem metryki.
\begin{equation}
	I[g]=-\kappa\int d^{d+1}x\sqrt{-g}\left(  \frac{1}{2^{p}}\ \delta_{\nu
		_{1}...\nu_{2p}}^{\mu_{1}...\mu_{2p}}\,R_{\mu_{1}\mu_{2}}^{\nu_{1}\nu_{2}%
	}\cdots R_{\mu_{2p-1}\mu_{2p}}^{\nu_{2p-1}\nu_{2p}}-2\Lambda\right)  \,.
	\label{PL}%
\end{equation}
Po uwzględnieniu stałej kosmologicznej teoria dopuszcza unikalne próżnie de Sittera (dS) lub anty-de Sittera (AdS). To sformułowanie grawitacji Pure Lovelock (PL) wykazuje podobieństwo do działania Einsteina-Hilberta, w szczególności posiadając unikalną próżnię o stałej krzywiźnie oraz możliwość rozwiązań w postaci naładowanych czarnych dziur.

Przeprowadzając analizę hamiltonowską tej teorii, możemy znaleźć pierwotne i wtórne więzy oraz warunki ich zachowania w czasie. Ostatecznie jesteśmy w stanie określić liczbę i charakter lokalnych symetrii oraz fizyczne stopnie swobody w teorii.

Analiza kanoniczna w $(d+1)$-wymiarowej czasoprzestrzeni ostatecznie wykazała zmienność liczby stopni swobody w teorii PL, w zakresie od $0$ do $\frac{(d+1)(d-2)}{2}$. Zakres ten jest zależny od tła, co stanowi rozróżnienie w porównaniu do standardowej Ogólnej Teorii Względności. Badanie potwierdza, że dla ogólnych przypadków tła, w tym przestrzeni (A)dS oraz statycznych czarnych dziur, teoria PL posiada taką samą liczbę stopni swobody jak OTW. Jednakże w innych tłach, takich jak przestrzeń płaska, teoria może stać się topologiczna i nie posiadać stopni swobody.

%%%%%%%%%%%%%%%%%%%%%%%%%%%%%%%%%%%%%%%%%%%%%%%%%%%%%%%%%%%%%%%%%%%%%%%%%%
\subsection{Podsumowanie}

Moje badania dotyczą rozwoju modeli grawitacji i leżących u ich podstaw symetrii, dostarczając nowatorskich struktur algebraicznych i wglądu w ich zastosowania zarówno w kontekście teoretycznym, jak i fizycznym. Badając algebry rezonansowe, rozszerzając ramy supersymetrii i stosując te odkrycia do praktycznych scenariuszy, takich jak izolatory topologiczne, supergrawitacja i termodynamika czarnych dziur, moja praca otwiera nowe możliwości dla dalszych badań w teoriach grawitacji i poza nimi. Innowacyjne metodologie opracowane w trakcie tych badań, w tym zaawansowane techniki obliczeniowe i algebra symboliczna, podkreślają rosnące znaczenie nowoczesnych narzędzi obliczeniowych w fizyce teoretycznej. 

W ramach swojej pracy badałem supersymetryczne rozszerzenia algebr rezonansowych, konstruując różne modele supergrawitacji w trzech i czterech wymiarach. Prace te obejmują identyfikację wszystkich możliwych rozszerzeń supersymetrycznych $\mathcal{N}=1$ i $\mathcal{N}=2$, jak również odpowiadających im działań. Moim najważniejszym wkładem w tym obszarze było wprowadzenie nowej klasy superalgebr Liego, wygenerowanych za pomocą metody obliczeniowej opartej na szukaniu wzorców, która dostarcza istotnych modyfikacji teorii cechowania.

Główną motywację stanowiła tu bezpośrednia inspiracja supersymetrią – rozszerzenie ogólne w celu uwzględnienia dodatkowych symetrii/pól oraz wyjście poza początkowe ramy Ogólnej Teorii Względności. Naturalnie uwaga wykracza poza konstruowanie struktur algebraicznych, koncentrując się również na ich zastosowaniu w budowie modeli grawitacyjnych. Godne uwagi przykłady obejmują teorie $\mathcal{N}=0$ (OTW), teorie $\mathcal{N}=1$ (standardowa supergrawitacja) oraz teorie $\mathcal{N}=2$ z dwoma superładunkami $Q_\alpha$ i $Y_\alpha$. Ramy algebr rezonansowych dodatkowo umożliwiają włączenie dodatkowych pól bozonowych i generatorów, oferując nowe sprzężenia w modelach grawitacyjnych i wzbogacając teoretyczny krajobraz.

Po odkryciu tak szerokiego spektrum przypadków konieczne staje się teraz uzasadnienie konkretnego wyboru algebry. Nowe przykłady mogą znacząco wzbogacić nasze zrozumienie powiązań między różnymi teoriami supergrawitacji, a także ujawnić unikalne, dotychczas niebadane właściwości. Struktury te mogą skrywać ukryte powiązania z takimi zjawiskami jak ciemna materia czy ciemna energia.

Chociaż moje badania koncentrowały się na strukturach opartych na AdS, opracowana tutaj metodologia może zostać rozszerzona na inne struktury algebraiczne. Zachowując stałe struktury i wprowadzając dodatkowe generatory, można uzyskać podobne uogólnienia \textit{rezonansowe}. Spekulacyjnie, podejście to może dostarczyć wglądu również w inne obszary, takie jak trzy rodziny cząstek w Modelu Standardowym.

Patrząc w przyszłość, planuję dalej badać zastosowanie algebr rezonansowych w teoriach bi-metrycznych oraz zbadać, w jaki sposób duże modele językowe (LLM) i techniki uczenia maszynowego (takie jak analiza semantyczna, analiza istotności i regresja symboliczna \cite{Udrescu:2019mnk}) mogą być zintegrowane z badaniami nad strukturami algebraicznymi.

%%%%%%%%%%%%%%%%%%%%%%%%%%%%%%% Bibliography %%%%%%%%%%%%%%%%%%%%%%%%%%%%%

\bibliography{Bibliografia}
\bibliographystyle{unsrt}
%\bibliographystyle{apsrev4-1}

%%%%%%%%%%%%%%%%%%%%%%%%%%%%%%%%%%%%%%%%%%%%%%%%%%%%%%%%%%%%%%%%%%%%%%%%%%%

\section{Prezentacja działalności naukowej i artystycznej}
%\section{Prezentacja znaczącej działalności naukowej lub artystycznej prowadzonej na więcej niż jednej uczelni, w instytucji naukowej lub kulturalnej, zwłaszcza zagranicznej.} 

\subsection*{Prezentacja działalności naukowej}

\indent\indent Od ukończenia pracy magisterskiej z fizyki teoretycznej w 2007 roku na Uniwersytecie Wrocławskim, moje badania koncentrują się na rozszerzonych modelach grawitacji oraz strukturach algebraicznych w fizyce.

W trakcie moich studiów doktoranckich na Uniwersytecie Wrocławskim, pod opieką prof. Jerzego Kowalskiego-Glikmana, badałem teoretyczne podstawy teorii grawitacji jako zdeformowanej teorii BF. Teoria ta obejmowała najbardziej ogólne działanie pierwszego rzędu, zawierające część Einsteina-Cartana i stałą kosmologiczną, a także człon Holsta oraz trzy niezmienniki topologiczne: Eulera, Pontryagina i Nieh-Yana. Badania te dotyczyły konstrukcji modeli grawitacyjnych i supergrawitacyjnych, analizy ich struktur algebraicznych oraz badania szerokiego spektrum fizycznych konsekwencji. Praca ta stanowiła solidną podstawę dla moich dalszych badań nad algebraicznymi rozszerzeniami teorii grawitacyjnych.

Po uzyskaniu stopnia doktora w 2012 roku otrzymałem grant FONDECYT na badania podoktorskie na Pontificia Universidad Católica de Valparaíso w Chile (2013-2016). W tym okresie, we współpracy z Olivierą Miskovic, Nelsonem Merino i Nareshem Dadhichem, badałem stopnie swobody w teorii grawitacji Lovelocka (\textbf{H1}). Nasze badania wykazały unikalne cechy dynamiczne modeli Pure Lovelock, ujawniając, że liczba fizycznych stopni swobody zależy od geometrii tła, co wyróżnia je w porównaniu ze standardową Ogólną Teorią Względności.

Chilijska społeczność badawcza, specjalizująca się w klasycznych modelach grawitacji i rozwiązaniach równań Einsteina, zapewniła mi inspirujące środowisko, które zachęciło mnie do podjęcia wyzwań, takich jak rozszerzenie algebr Poincar\'e i AdS. Praca ta doprowadziła do znaczących odkryć dotyczących rozszerzonych algebr (tzw. S-expanded), a później algebr rezonansowych, które odegrały kluczową rolę w konstruowaniu nowych modeli grawitacyjnych (\textbf{H2, H3}). Współpraca z Nelsonem Merino, Patrickiem Conchą i Evelyn Rodríguez zaowocowała później kilkoma kluczowymi publikacjami, które pogłębiły moje zrozumienie algebraicznych rozszerzeń w grawitacji.

W 2017 roku rozpocząłem pracę na stanowisku adiunkta w Instytucie Fizyki Teoretycznej Uniwersytetu Wrocławskiego, gdzie moje badania doprowadziły do kilku osiągnięć: zastosowań algebr rezonansowych i ich roli w opisie izolatorów topologicznych (\textbf{H4}) oraz charakteryzacji wszystkich możliwych algebr podobnych do algebry Maxwella (\textbf{H5}) dla większej liczby generatorów. Po powrocie do Polski wznowiłem współpracę z prof. J. Kowalskim-Glikmanem. Nasze wspólne projekty bardzo skorzystały na moim doświadczeniu podoktorskim, szczególnie w obszarach takich jak symetrie asymptotyczne i człony topologiczne w grawitacji. Uzyskaliśmy nowe spojrzenie na struktury brzegowe teorii grawitacyjnych, czego kulminacją były publikacje na temat teorii BF (\textbf{H7}) oraz badania porównawcze różnych sformułowań grawitacji dla przestrzeni asymptotycznie AdS (\textbf{H11}).

W 2020 roku zostałem mianowany na stałe stanowisko adiunkta w Instytucie Fizyki Teoretycznej Uniwersytetu Wrocławskiego, gdzie obecnie kontynuuję swoje badania. Od tego momentu moja uwaga zaczęła skupiać się na symbolicznych obliczeniach komputerowych oraz innych nowoczesnych technologiach stosowanych w badaniach teoretycznych. Moja współpraca z Krzysztofem Graczykiem, specjalizującym się w ostatnich latach w uczeniu maszynowym w fizyce, doprowadziła do kilku publikacji dotyczących rezonansowych superalgebr oraz ich zastosowań w modelach supergrawitacji (\textbf{H9, H10, H12}). Badania te rozszerzyły krajobraz możliwych teorii supergrawitacji poprzez identyfikację nowych struktur algebraicznych, które uwzględniają dodatkowe generatory bozonowe i fermionowe.

Moje badania konsekwentnie dotyczą struktur algebraicznych, modeli grawitacyjnych i technik obliczeniowych, przyczyniając się do głębszego zrozumienia podstawowych symetrii w grawitacji wraz z ich zastosowaniami. Jestem zaangażowany w dalszy rozwój tej dziedziny poprzez kontynuację współpracy, innowacyjne badania i wykorzystanie nowoczesnych metod obliczeniowych.

%%%%%%%%%%%%%%%%%%%%%%%%%%%%%%%%%%%%%%%%%%%%%%%%%%%%%%%%%%%%%%%%%%%%%%%%%%%%%%%%
\subsection*{Prezentacja działalności artystycznej}

\indent\indent Byłem jednym z głównych współtwórców  i realizatorów projektu edukacyjno - rekreacyjnego parku „Bulwar Fizyków”. Efektem projektu było powstanie instalacji Zwierciadeł Akustycznych umieszczonych nad rzeką Odrą obok Wydziału Fizyki i Astronomii, sfinansowanej z budżetu miejskiego Wrocławia. Instalacja stała się jednym z charakterystycznych punktów miasta, prezentującym fizykę klasyczną studentom i turystom.

Pełniłem rolę konsultanta naukowego przy produkcji słuchowiska sf „Fantom” autorstwa polskiego pisarza Marcina Ciszewskiego, wyprodukowanej przez Audiotekę.

Ponadto, jako pierwszy odnalazłem akt urodzenia Stefana Banacha w archiwach Krakowa. Moje badania nad skomplikowanymi korzeniami największego polskiego matematyka zostały opublikowane w polskim czasopiśmie genealogicznym „More Maiorum” (2017).

%%%%%%%%%%%%%%%%%%%%%%%%%%%%%%%%%%%%%%%%%%%%%%%%%%%%%%%%%%%%%%%%%%%%%%%%%%%%%%%%
\section{Dydaktyka, organizacja i popularyzacja nauki}
% 6. Prezentacja osiągnięć dydaktycznych, organizacyjnych oraz w zakresie popularyzacji nauki lub sztuki

\subsection{Doświadczenie dydaktyczne}

Na przestrzeni lat prowadziłem szeroki zakres wykładów, zajęć laboratoryjnych i ćwiczeń, obejmujących różnorodne tematy i zapewniających kompleksową edukację studentom na różnych poziomach. Prowadzone wykłady obejmują m.in. \textit{Praktyczną Mechanikę Kwantową}: kurs łączący wykłady z zajęciami laboratoryjnymi, kładący nacisk na wykorzystanie oprogramowania Mathematica do obliczeń i wizualizacji w mechanice kwantowej, co wzmacnia praktyczne zrozumienie i zastosowania. \textit{Matematykę dla ISSP}: kurs matematyki dla pierwszego roku, dostosowany do specjalizacji Informatyka i Systemy Pomiarowe. \textit{Obliczenia Numeryczne i Symboliczne w Fizyce}: zaawansowane zajęcia laboratoryjne koncentrujące się na wykorzystaniu Mathematici i języka Wolfram do modelowania, symulacji i analizy danych w różnych zastosowaniach fizycznych. \textit{Astrophysics 2}: będące kontynuacją głównego kursu astrofizyki.\\
%
\\
Pełna lista prowadzonych zajęć:
\begin{itemize}
	\item \textbf{Wykłady:} Astrophysics 2: \textsl{2024}
	\item \textbf{Wykłady:}  Physics I, Physics II w ramach "Programu NAWA \textsl{2022-2023}" pod patronatem Politechniki Wrocławskiej (dwu-semestralny kurs fizyki dla obcokrajowców z Kartą Polaka)
	\item \textbf{Wykłady:} Matematyka dla ISSP II: \textsl{2024, 2023, 2022, 2021, 2020, 2019}
	\item \textbf{Wykłady:} Matematyka dla ISSP I: \textsl{2024, 2023, 2022, 2021, 2020, 2019, 2018}
	\item \textbf{Wykłady:} Praktyczna Mechanika Kwantowa: \textsl{2023, 2022, 2021}
	\item Advance Cosmology (ćwiczenia): \textsl{2024}
	\item Laboratorium komputerowe metod matematycznych (laboratorium): \textsl{2024}
	\item Praktyczna Mechanika Kwantowa (ćwiczenia): \textsl{2023, 2022, 2021, 2020, 2019, 2018}
	\item Matematyka dla ISSP II (ćwiczenia): \textsl{2023, 2022, 2021, 2020, 2019}
	\item Obliczenia numeryczne i symboliczne w fizyce (laboratorium): \textsl{2023, 2022, 2021, 2020, 2019, 2018}
	\item Programy użytkowe (laboratoria): \textsl{2019, 2018, 2017}
	\item Matematyka dla Informatyków II (laboratorium): \textsl{2019}
	\item Matematyka dla Informatyków (2 laboratoria): \textsl{2018}
	\item Algorytmy i struktury danych (ćwiczenia): \textsl{2018}
	\item Algebra II (ćwiczenia): \textsl{2018}
	\item Matematyka dla Informatyków (laboratorium): \textsl{2017}
	\item Algebra I (ćwiczenia): \textsl{2017}
	\item Technologie komputerowe (5 grup): \textsl{2013}
	\item Mechanika Teoretyczna (ćwiczenia): \textsl{2012}
	\item Programy biurowe (ćwiczenia): \textsl{2012}
	\item Programy użytkowe (ćwiczenia): \textsl{2009}
	\item Elektrodynamika (ćwiczenia): \textsl{2008}
\end{itemize}

%%%%%%%%%%%%%%%%%%%%%%%%%%%%%%%%%%%%%%%%%%%%%%%%%%%%%%%%%%%%%%%%%%%%
\subsection{Wykłady i Seminaria Popularyzatorskie}

\begin{itemize}
	\item 2023: Laboratorium AI dla Fizyki (LAIP) (instytutowa grupa naukowa) \textsl{"AI Feynman: A Physics-Inspired Method for symbolic Regression"}
	\item 2023: $25$-lecie obserwatorium astronomicznego w Głogowie: \textsl{Od Ziemi do krańców Wszechświata}
	\item 2022: wykład na Dolnośląskim Festiwalu Nauki we Wrocławiu, \textsl{Kostka Rubika i wyższe wymiary}
	\item 2021: wykład na otwarciu roku akademickiego, WFiA \textsl{Witamy na Fizyce}
	\item 2018: wykład w ramach Wykładów Fizycznych na Uniwersytecie Wrocławskim
	\item 2018: wykład na Dniach Otwartych Uniwersytetu Wrocławskiego
	\item 2011: wykład na Dolnośląskim Festiwalu Nauki we Wrocławiu, \textsl{Fizyka teoretyczna – gdzie, kiedy i jak?}
	\item 2009: wykład dla gimnazjum we Wrocławiu, \textsl{Skąd to się wzięło?}
	\item 2008-2012: liczne seminaria doktoranckie, 2006-2008: liczne seminaria magisterskie, wykłady dla koła naukowego: "Kosmologia", "Kostki Rubika w wyższych wymiarach", "Stanisław Ulam: historia matematyka".
\end{itemize}

%%%%%%%%%%%%%%%%%%%%%%%%%%%%%%%%%%%%%%%%%%%%%%%%%%%%%%%%%%%%%%%%%%%%
\subsection{Organizacyjne: Konferencje i Wydarzenia}

\begin{itemize}
	\item Współorganizator 59. Zimowej Szkoły Fizyki Teoretycznej, 9-22 lutego 2023, Wojanów, Polska
	\item Współorganizator Sympozjum na temat kwantowych symetrii, niekomutatywnej geometrii i kwantowej grawitacji, 24-25 czerwca 2021, Wrocław, Polska
	\item Współorganizator 5. Konferencji Polskiego Towarzystwa Relatywistycznego, 24-27 września 2018, Wojanów, Polska
	\item Organizator 54. Zimowej Szkoły Fizyki Teoretycznej, 18-24 lutego 2018, Lądek-Zdrój, Polska
	\item Organizator "Chilean Gravity Meeting in Buenos Aires", 20-21 listopada 2014, Argentyna
	\item Współorganizator pierwszego "Meeting on the Horizon", 10-14 marca 2014, Valparaíso, Chile
	\item Sekretarz i współorganizator konferencji "Planck Scale 2009" we Wrocławiu, Polska
\end{itemize}

%%%%%%%%%%%%%%%%%%%%%%%%%%%%%%%%%%%%%%%%%%%%%%%%%%%%%%%%%%%%%%%%%%%%
\subsection{Aktywność i obowiązki na Wydziale i poza nim}

\begin{itemize}
	\item Członek Wydziałowego Zespołu ds. Oceny Jakości Kształcenia w 2024 roku
	\item Wydziałowy Koordynator ds. nauczania online od 2020 roku
	\item Członek i webmaster Polskiego Towarzystwa Relatywistycznego od 2019 roku
	\item Przedstawiciel ds. relacji publicznych na wydziale (2020-2021)
	\item Organizator Kina Naukowego: \textit{Universyteka} (pierwsza edycja: grudzień 2017, druga edycja: kwiecień 2018)
	\item Silne zaangażowanie w inicjatywy związane z budową nowego budynku dla Wydziału Fizyki i Astronomii
	\item Organizator Wrocław Rubik Fewest Moves Challenge 2019
	\item Współorganizator $5$-tych Mistrzostw Europy w Kostce Rubika, Wrocław 2012 (z udziałem 300 zawodników)
	\item Wielokrotne wyróżnienia Nagrodą Rektora Uniwersytetu Wrocławskiego za działalność organizacyjną w Instytucie Fizyki Teoretycznej
\end{itemize}

%%%%%%%%%%%%%%%%%%%%%%%%%%%%%%%%%%%%%%%%%%%%%%%%%%%%%%%%%%%%%%%%%%%%
\section{Dodatkowe informacje}
% 7. Apart from the information set out in 1-6 above, the applicant may include other information about his/her professional career, which he/she deems important.  

%%%%%%%%%%%%%%%%%%%%%%%%%%%%%%%%%%%%%%%%%%%%%%%%%%%%%%%%%%%%%%%%%%%%

\subsection*{Opieka nad pracami dyplomowymi studentów}

Byłem opiekunem kilku studentów realizujących prace magisterskie, licencjackie oraz inżynierskie, koncentrujące się na tematach od paradoksu informacyjnego po obliczenia symboliczne w fizyce oraz nowoczesne technologie:
\begin{itemize}
	\item Praca magisterska: Michał Bobula (2021): "Information Paradox: from Hawking Radiation to Page Curve"
	\item Praca magisterska: Mradul Navhal (2018): "Mathematica based calculations of the Kretschmann scalar and other terms like Gauss-Bonnet, Pontryagin, and Holst"
	\item Praca licencjacka: Michał Popiel (2022): "Wykorzystanie Wolfram Mathematica do rozwiązywania równań Einsteina"
	\item Praca licencjacka: Gustaw Tański (2020): "Symboliczne wizualizacje koneksji Christoffela z wykorzystaniem silnika Unity"
	\item Prace inżynierskie: D. Grafik ('23), K. Sobczyk ('22), B. Nawrocki ('22), R. Hupka ('22), B. Urban ('21)
\end{itemize}

%%%%%%%%%%%%%%%%%%%%%%%%%%%%%%%%%%%%%%%%%%%%%%%%%%%%%%%%%%%%%%%%%%%%

\subsection*{Udział w inicjatywach międzynarodowych}

\begin{itemize}
	\item Byłem członkiem akcji COST CA18108 (QG-MM) „Quantum gravity phenomenology in the multi-messenger approach”.
	\item Jestem również członkiem Polskiego Towarzystwa Relatywistycznego (i osobą odpowiedzialną za ich stronę internetową).
	\item W latach 2018 i 2019 uczestniczyłem w programie wymiany pedagogicznej ERASMUS+ (Czechy, Włochy).
	\item Od początku studiów doktoranckich i przez cały okres mojej pracy badawczej regularnie uczestniczyłem, wygłaszając referaty na międzynarodowych konferencjach i warsztatach. Prowadziłem również seminaria instytutowe na różnych uniwersytetach i w instytutach. Pełna lista znajduje się w załączonym do dokumentów „Wykazie osiągnięć naukowych”.
\end{itemize}

%%%%%%%%%%%%%%%%%%%%%%%%%%%%%%%%%%%%%%%%%%%%%%%%%%%%%%%%%%%%%%%%%%%%
\subsection*{Stypendia i Granty}

W trakcie mojej kariery akademickiej, zarówno przed, jak i w trakcie realizacji doktoratu, a także na etapie podoktorskim, zrealizowałem kilka grantów:
\begin{itemize}
	\item 2011-2012: Byłem kierownikiem grantu promotorskiego „Grawitacja, termodynamika i topologiczne teorie pola”, projektu wspieranego przez grant numer N N 202 112 740 z Narodowego Centrum Nauki (NCN).
	\item 2011-2013: Byłem kierownikiem grantu PRELUDIUM dla młodych naukowców z NCN, „Kanoniczne i termodynamiczne aspekty grawitacji jako teorii BF z więzami”.
	\item 2013-2016: Otrzymałem stypendium podoktorskie FONDECYT, przyznane przez Chilijską Komisję ds. Badań Naukowych i Technologicznych (Comisión Nacional de Investigación Científica y Tecnológica), realizowane od listopada 2013 do listopada 2016. Dzięki temu pracowałem w PUCV w Valparaíso w Chile, co umożliwiło mi współpracę z różnorodną społecznością naukową i zdobycie bezcennych doświadczeń w mojej dziedzinie badań.
\end{itemize}

%%%%%%%%%%%%%%%%%%%%%%%%%%%%%%%%%%%%%%%%%%%%%%%%%%%%%%%%%%%%%%%%%%%%

%%%%%%%%%%%%%%%%%%%%%%%%%%%%%%%%%%%%%%%%%%%%%%%%%%%%%%%%%%%%%%%%%%%%
~~\\[40pt]
\noindent
\begin{minipage}{\textwidth}
	\hfill % Wyrównanie całości do prawej
	\begin{minipage}{0.4\textwidth} % Ustaw szerokość dla obrazka
		\includegraphics[width=0.8\textwidth]{podpis_a.PNG} % Wstaw obrazek
		\\[-10pt]
		\rightline{.......................................................................} 
		\rightline{(Podpis wnioskodawcy)\qquad\qquad\qquad}
	\end{minipage}
\end{minipage}
\end{document}